Attention and Transformers Are Gauge-Theoretic
Variational Free Energy Minimization
Robert C. Dennis cdenn016@gmail.com
Independent Researcher
Leander, Texas 78641, USA
Abstract
Transformer attention is variational inference over information sources. Modeling each
token as a Gaussian agent on a statistical fiber bundle, we derive the generalized attention
weight βij = softmax(−DKL[qi∥Ωij qj ]/τ ) from a single variational principle: minimizing
the free energy of a mixture-of-sources generative model in which each agent infers, via
gauge transport, which neighbor generated its state. The KL divergence arises exactly;
the softmax follows from constrained optimization over the categorical source-selection
posterior.
From this principle, transformer architectural choices emerge as consequences of the
variational geometry: the temperature scaling 1/√dk from dimensional concentration of
the KL divergence, layer normalization as the geometric condition for frame-independent
inference, multi-head attention as a block-diagonal restriction of the gauge algebra, and
causal masking and positional biases from non-uniform attention priors πj . Two succes-
sive limits (isotropic covariances, flat gauge connection) recover the standard rule βij ∝
softmax(QiK⊤
j /√dk), while the GL(K) invariance of KL divergence implies that WQ, WK
are themselves gauge transformations with WQW ⊤
K = σ−2Ω−⊤ parametrizing a learned
gauge transport, suggesting that standard transformers already implement gauge-theoretic
attention. Gradient descent on the free energy provides a variational interpretation of the
transformer training loss, with the chain rule through layered compositions corresponding
to variational message-passing.
We validate the unreduced theory by training a GL(K) gauge transformer on WikiText-
103 using only KL divergences, gauge transport, and natural gradient dynamics (without
learned attention projections, MLPs, or pointwise activation functions such as ReLU or
GELU). The model achieves test perplexity 74.9, surpassing a matched Kneser-Ney 5-
gram baseline (PPL 134.8) by 1.80×, with emergent semantic structure in the learned
gauge frames. We separately validate the degenerate-limit predictions on frozen BERT
across 105 diverse passages, confirming quantitative agreement between KL-based and dot-
product attention (grand mean r = 0.804 at the predicted optimal temperature, 95% CI
[0.771, 0.838], with cross-passage SE < 0.02 for all 144 heads). Both architectures operate in
the flat bundle limit with trivially vanishing holonomy: a geometric regime we conjecture is
matched to the approximately path-independent semantic transport that language requires.
Keywords: gauge theory, free energy principle, transformer attention, variational infer-
ence, information geometry, natural gradient, symmetry breaking, multi-agent systems,
GL(K) invariance
1 Introduction
Language is fundamentally a dynamic informational system: speakers and writers encode
beliefs under uncertainty, listeners decode them, and the statistical regularities of this pro-
©2026 Robert C. Dennis.
License: CC-BY 4.0, see https://creativecommons.org/licenses/by/4.0/.
Dennis
cess reflect the structure of communication itself. Language modeling is a snapshot of this
dynamic system, and the architectures that excel at it should reflect the mathematical
structure of information exchange under uncertainty.
Recent advances in neuroscience and intelligent systems have independently converged
on the idea that intelligent systems integrate perception, inference, and communication un-
der the constraints of uncertainty (Bahdanau et al., 2014; Amari, 1998; Bronstein et al.,
2021; Foerster et al., 2016; Wooldridge, 2009). Friston’s Free Energy Principle (FEP) pro-
vides a general variational formulation of inference in cognitive systems (Friston, 2010; Parr
et al., 2022; Friston et al., 2017; Ramstead et al., 2020), whereas the attention mechanism
in modern machine learning architectures defines a powerful (although empirically derived)
rule for token prediction (Clark et al., 2019; Vaswani et al., 2017). Despite their shared
reliance on probabilistic inference and pairwise interaction, these two frameworks remain
stubbornly separated. In particular, transformer attention lacks an underlying geometric
or mathematical foundation and details on how and why modern machine learning archi-
tectures operate as well as they do remains obscure.
Many varieties of transformer and attention architectures have been proposed, imple-
mented, and studied in recent years (Fuchs et al., 2020; Thomas et al., 2018). Some ar-
chitectures make use of bundle geometric frameworks but lack a first-principles foundation
or connection to the FEP (Kondor and Trivedi, 2018; Finzi et al., 2020; Bronstein et al.,
2021). Furthermore, attempts at curved space token embeddings have produced mixed
results (Bonnabel, 2013; Absil et al., 2008).
Several complementary mathematical perspectives on attention have been developed.
Tsai et al. (2019) and Katharopoulos et al. (2020) interpret attention through the lens of
kernel methods, showing that softmax attention approximates a kernel function and can
be linearized. Ramsauer et al. (2021) demonstrate that attention implements the update
rule of modern Hopfield networks, providing an energy-based perspective. The information
bottleneck principle (Tishby and Zaslavsky, 2015; Shwartz-Ziv and Tishby, 2017) offers
a related information-theoretic view of layer-wise representation learning. On the side of
backpropagation-free learning, Hinton (2022) proposes the forward-forward algorithm as an
alternative to backpropagation. In neuroscience, predictive coding (Rao and Ballard, 1999;
Bogacz, 2017; Millidge et al., 2021) implements variational free energy minimization through
local prediction errors, sharing the variational foundation of our framework but lacking
the gauge-theoretic structure that connects to attention mechanisms. Our work differs
from these approaches in providing a unified geometric framework from which attention,
temperature scaling, layer normalization, and training dynamics all emerge as consequences
of a single variational principle on a statistical fiber bundle.
In this report we propose a first-principles, gauge-equivariant framework that connects
these domains. Our framework is based on a principal bundle geometry whereby each agent
is modeled as a smooth local section of an associated bundle with statistical manifold fibers
over a ”noumenal” base manifold. Inter-agent communication arises naturally through a
(generally) non-abelian gauge connection that defines parallel-transport operators between
agents’ local gauge frames. Within this geometry, attention emerges as a gauge-aligned
Kullback–Leibler (KL) term derived directly from the variational free energy of a coupled
multi-agent generative model based on mixture Gaussians.
2
Attention and Transformers as Free Energy Minimization
We then show that in the flat-bundle, isotropic, delta-function limit attention reduces
to the standard transformer dot-product attention QKT , thereby identifying attention as
a degenerate case of a broader geometric and statistical law of communication predicated
upon the FEP and information geometry. We demonstrate this by training a transformer
based upon variational gradient descent of a generalized free energy of multivariate Gaussian
agents as well as by applying our gauge attention theory to a frozen transformer (BERT)
across a corpus of 105 diverse passages, with cross-passage validation confirming the ro-
bustness of our predictions. We then describe numerical results that suggest that token
encoding under data training (i.e. agent/token belief alignment) is a symmetry breaking
phenomenon of our gauge-theoretic framework.
Finally, we show that gradient descent on the gauge-equivariant free energy, in the de-
terministic constant-gauge limit, recovers the standard transformer training update. This
provides a variational interpretation of the training loss: the free energy functional maps
onto the loss plus regularization, and the chain rule through layered compositions cor-
responds to variational message-passing. We emphasize that the chain rule is a general
property of differentiable function composition; what the framework contributes is an in-
terpretation of the loss function and its gradients as variational quantities, not a derivation
of backpropagation itself. The framework naturally invokes natural gradient descent on the
Fisher-Rao metric, which offers theoretical convergence advantages; whether this translates
to practical speedup in deep learning requires controlled empirical comparison, which we
leave to future work.
2 Methods
We begin with a general geometric construction that naturally supports hierarchical emer-
gence of meta-agents and scale interactions. The details of the most general geometry and
framework can be found in the appendix and references (Nakahara, 2003; Kullback and
Leibler, 1951; Sternberg, 1994; Fulton and Harris, 1991; Frankel, 2011; Baez and Muniain,
1994).
Briefly, in this current report, we describe agents as modeled by a local section of
an associated bundle to a principal G fiber bundle whereby tokens are represented by
agents occupying a single 0-dimensional base manifold. Although the base manifold is 0-
dimensional and carries no intrinsic ordering or metric, the discrete index set of agents
(tokens) may be endowed with additional structure such as a sequential ordering or a
distance function that enters the theory only through the attention prior πj (Section 3.4.5)
and/or through position-dependent gauge frames (Section 4.5). This separation between the
geometric base (which determines gauge connections and curvature) and the combinatorial
index structure (which determines masking and positional biases) is important: causal
masking and distance-dependent priors impose structure on which agents communicate,
not on the underlying fiber geometry. We shall show that this framework enables a unified
description of both intra-agent inference and inter-agent communication in terms of gauge
frame transport and alignment via a generalized free energy principle.
Notation and glossary. Table 1 summarizes the principal symbols used throughout this
work.
3
Dennis
Base manifold C
Ui
Uj
Ui ∩ Uj
Bq ∼= RK
qi(c)
qj (c)
c∗
ϕi
ϕj
Ωij qj
Ωij (c∗)
DKL(qi∥Ωij qj )
Figure 1: Agent belief sections over the associated bundle Eq = N ×ρq Bq. The shaded
surface is the base manifold C; vertical dashed lines are belief fibers Bq ∼= RK . Agents i
and j each have a local open neighborhood on C: Ui (blue, dashed) and Uj (red, dashed).
Each agent’s section is defined only over its own neighborhood: the blue section qi(c) exists
over Ui, the red section qj (c) over Uj . In the overlap Ui ∩ Uj , both sections coexist in
the same fibers. At the highlighted fiber over c∗ ∈ Ui ∩ Uj , qi(c∗) and qj (c∗) reside in
the same fiber but are expressed in different local gauge frames ϕi, ϕj (purple axes). The
transport operator Ωij (c∗) = exp(ϕi) exp(−ϕj ) (green arrow) acts within this fiber, rotating
qj from j’s frame into i’s frame to produce Ωij qj (dotted red ellipse). The attention cost
DKL(qi∥Ωij qj ) is then the divergence between same-frame distributions within a single fiber.
For transformers (dim C = 0), all tokens share one base point, the neighborhoods collapse
to a single point, and the entire computation occurs within one fiber.
2.1 An intuitive simplification for non-geometers
For interdisciplinary readers, it may be helpful to visualize each agent in this framework as
a vector and a matrix field defined over a finite spatial region.
Concretely, at every point c in the base manifold, agent i carries a mean vector μi(c) and
a covariance matrix Σi(c) representing, respectively, its local expected state and uncertainty.
Together these form a smooth Gaussian field—the agent’s belief section of the statistical
bundle.
We will see that the agent’s local gauge frame ϕi(c) may be intuitively interpreted as
the agent’s subjective frame of reference; its internal coordinate system through which
beliefs and models are represented and compared. In this sense, ϕi(c) captures an agent’s
“internal orientation” toward the world: it determines how external information is projected
into the agent’s internal representational space. Gauge transformations ϕi(c) → ϕi(c) + ξ(c)
4
Attention and Transformers as Free Energy Minimization
thus correspond to changes in perspective that leave the underlying informational content
invariant, much like shifts of viewpoint in perception that preserve the same world model.
2.1.1 Gauge Group and Representations
The natural gauge group for KL-based attention is G = GL(K), the general linear group
of invertible K × K matrices. This acts on the fibers via representations:
ρq : GL(K) → GL(dq), ρp : GL(K) → GL(dp). (1)
For Ω ∈ GL(K), the gauge action on a Gaussian state is:
ρq(Ω) · (μq, Σq) = ρq(Ω) μq, ρq(Ω) Σq ρq(Ω)⊤,
ρp(Ω) · (μp, Σp) = ρp(Ω) μp, ρp(Ω) Σp ρp(Ω)⊤. (2)
By Theorem 1, KL divergence is invariant under this action for any Ω ∈ GL(K). The
only requirement is invertibility (det Ω̸ = 0), not orthogonality. (Since our transport oper-
ators are parameterized via the matrix exponential, det(exp(ϕ)) = exp(tr(ϕ)) > 0, so their
determinants are always positive.)
2.1.2 Local Gauge Frames
Each agent i has a gauge frame field ϕi : Ui → gl(K) which may, in general, vary across
the base manifold. On higher-dimensional base manifolds, a gauge frame field induces a
local connection one-form A(i)
μ (c) = U −1
i (c) ∂μUi(c) (where Ui = exp(ϕi)) and associated
field strength F (i)
μν = ∂μA(i)
ν − ∂ν A(i)
μ + [A(i)
μ , A(i)
ν ]. However, a connection derived from a
single-valued gauge function ϕi(c) in this way is pure gauge and has identically vanishing
curvature Fμν = 0. Non-trivial curvature and holonomy require independent connection
variables on edges or paths, transforming as Ωij 7 → giΩij g−1
j rather than being determined
by vertex-local group elements (see Appendix A for details).
In the 0-dimensional case, which we consider here, there exist no gauge connections
Aμ or field strengths Fμν . What remains is the algebraic structure of group-valued frame
transformations Ωij = exp(ϕi) exp(−ϕj ) acting on statistical fibers. Since these trans-
ports are constructed from vertex-local group elements, they satisfy the cocycle condition
Ωij Ωjk = Ωik and have trivially vanishing discrete holonomy Ωij ΩjkΩki = I for every closed
loop (Section 6.3). This algebraic structure already suffices to recover transformer atten-
tion (as shown below), and already provides a functional language model when implemented
without taking the degenerate limits (Section 5.3). Architectures with non-trivial holonomy
would require promoting the transport Ωij to an independent edge variable; this represents
a direction for future work (see Section 6.3).
2.1.3 GL(K) Gauge Invariance of KL Divergence
The Kullback-Leibler divergence is invariant under simultaneous invertible linear transfor-
mations of both distributions. Specifically, the KL divergence (and indeed all f -divergences)
is invariant under any common bi-jective change of variables (since the Jacobian factors can-
cel in the density ratio). By restricting to transformations that preserve the MVG family
5
Dennis
and act linearly on the sufficient statistics (μ, Σ) we find GL(K) as the symmetry group.
This is a standard result in information geometry (see, e.g., Amari (2016)), but its impli-
cations for transformer attention have not been previously recognized. We state it here for
completeness.
Theorem 1 (GL(K) Gauge Invariance) For any two Gaussian distributions P = N (μP , ΣP )
and Q = N (μQ, ΣQ) on RK , and any invertible linear transformation Ω ∈ GL(K), the KL
divergence is invariant under simultaneous push-forward:
DKL(Ω∗P ∥ Ω∗Q) = DKL(P ∥ Q), (3)
where Ω∗P = N (ΩμP , ΩΣP Ω⊤) denotes the pushforward of P under Ω.
Proof The KL divergence between Gaussians consists of three terms:
DKL(P ∥Q) = 1
2

Tr(Σ−1
Q ΣP ) + (μP − μQ)⊤Σ−1
Q (μP − μQ) − K + log det ΣQ
det ΣP

. (4)
Under the pushforward Ω∗P = N (ΩμP , ΩΣP Ω⊤) and Ω∗Q = N (ΩμQ, ΩΣQΩ⊤):
Trace term:
Tr(ΩΣQΩ⊤)−1(ΩΣP Ω⊤) = TrΩ−⊤Σ−1
Q Ω−1ΩΣP Ω⊤ = Tr(Σ−1
Q ΣP ). (5)
Quadratic term:
(ΩμP − ΩμQ)⊤(ΩΣQΩ⊤)−1(ΩμP − ΩμQ)
= (μP − μQ)⊤Ω⊤Ω−⊤Σ−1
Q Ω−1Ω(μP − μQ) = (μP − μQ)⊤Σ−1
Q (μP − μQ). (6)
Log-determinant term:
log det(ΩΣQΩ⊤)
det(ΩΣP Ω⊤) = log (det Ω)2 det ΣQ
(det Ω)2 det ΣP
= log det ΣQ
det ΣP
. (7)
The (det Ω)2 factors cancel identically. Since all three components are invariant,
DKL(Ω∗P ∥Ω∗Q) = DKL(P ∥Q).
Importantly, the proof extends to all f -divergences: i.e. for any convex function f such
that f (1) = 0,
Df (Ω∗P ∥Ω∗Q) = Df (P ∥Q), ∀ Ω ∈ GL(K), (8)
since the factors | det Ω| from the density transformation cancel in the ratio p(x)/q(x).
From a practical point of view, this implies that transport operators need only satisfy
det(Ωij )̸ = 0 (invertibility). Therefore, the full GL(K) acts as a gauge symmetry. It is
important to point out that our exponential parameterization eϕ restricts to det > 0 (since
det(exp(ϕ)) = exp(tr(ϕ)) > 0), but this is merely the identity component. However, this
choice is standard for gauge connections and does not present any imposition in our proof-
of-principle experiments. Disconnected transformations are reserved for future study.
Finally, we shall demonstrate that the query and key projection matrices WQ, WK ∈
GL(dk) in standard transformers are themselves related to gauge transformations. The
constant gauge limit corresponds not to the absence of gauge structure, but to a constant
gauge transformation satisfying WQW ⊤
K = σ−2Ω−⊤ shared among all token pairs.
6
Attention and Transformers as Free Energy Minimization
Notation conventions. We use K (or Kq) for the belief fiber dimension (embedding di-
mension), which generally does not coincide with the representation dimension of the gauge
group. For example, we may embed in 10 dimensions with two copies of the fundamental
GL(5) group. When discussing the reduction to standard transformers, dk denotes the per-
head dimension (dk = K/H for H heads). We use qi to denote both the distribution and
its density qi(ki), with the meaning clear from context. In what follows we will take the
gauge group dimension to be equal to the embedding dimension without loss of generality.
3 Gauge-Covariant Variational Free Energy
We now construct the variational free energy for a multi-agent system coupled by gauge
transport. All probability distributions are defined at a specific base manifold point c = c∗
unless otherwise noted. We label the fiber’s latent coordinates as ki ∈ B in order to define
the necessary integrals.
3.1 Agent State Space
Each agent i maintains two distinct latent variables living in separate fiber bundles:
ki ∈ RKq (belief latent in Eq), mi ∈ RKp (model latent in Ep). (9)
The full state of agent i is taken to be a multi-variate Gaussian (MVG) distribution
over each latent:
qi(ki) = N (ki ; μq,i, Σq,i), Σq,i ≻ 0,
si(mi) = N (mi ; μp,i, Σp,i), Σp,i ≻ 0. (10)
Each latent has an independent Gaussian base prior encoding agent-specific inductive
biases:
pi(ki) = N (ki ; μ(q)
0,i , Σ(q)
0,i ), ri(mi) = N (mi ; μ(p)
0,i , Σ(p)
0,i ). (11)
These priors are local in that they live within each agent’s own gauge frame ϕ and need
not be related across agents until transported via the gauge connection.
3.2 Gauge Transport
The gauge transport operators Ωij , ˜Ωij are pointwise gauge frame transformations:
Ωij (c) = eϕi(c) · e−ϕj (c) ∈ GL+(K), (12)
where ϕi : Ui → gl(K) is agent i’s gauge frame field (an open local section of the Lie
algebra bundle) and GL+(K) = {A ∈ GL(K) : det A > 0} is the identity component of
GL(K). These act on the latent variables as
Ωij kj := ρq(Ωij ) kj ∈ RKq ,
˜Ωij mj := ρp( ˜Ωij ) mj ∈ RKp , (13)
mapping agent j’s latent into agent i’s frame.
7
Dennis
3.3 Single-Agent Free Energy Principle
Before introducing multi-agent communication, we recall the standard variational free en-
ergy for a single agent i (Friston, 2010; Parr et al., 2022). Agent i maintains a belief qi(ki)
over its latent state, a prior pi(ki) encoding expectations, and receives observations oi. The
free energy principle posits that the agent minimizes:
Fsingle
i = DKL(qi∥pi) − Eqi [log p(oi | ki)]. (14)
The first term regularizes beliefs toward the prior and the second couples beliefs to
observations. An analogous expression holds for a model channel (Appendix F.1). We now
answer the question: how do agents communicate?
3.4 Deriving Attention from a Mixture-of-Sources Model
We will derive the agent-agent interaction terms and their softmax attention weights from
first principles by formulating each agent’s inference as a source selection over their gauge-
transported neighbors.
3.4.1 Mixture-of-Sources Generative Model
Agent i (the query) seeks to update its belief qi(k) by attending to a set of source agents
j (the keys). Agent i’s generative model assumes its latent state k was drawn from one of
the neighboring agents, selected by a categorical latent variable z ∈ {1, . . . , N }:
P (k, z) = P (k | z) P (z), (15)
where P (z = j) = πj is the prior probability of attending to agent j (the attention
prior), and P (k | z = j) = N (k; Ωij μj , Ωij Σj Ω⊤
ij ) is agent j’s belief transported into agent
i’s gauge frame via Ωij .
This construction leads to a generative model consistent with our required gauge struc-
ture. Agent j’s belief qj (kj ) lives in j’s local frame, and Ωij qj (kj ) is agent i’s representation
of what agent j believes, as seen by agent i. The categorical variable z endows the model
with probabilistic source selection whereby attention is considered as posterior inference
over which neighbor generated agent i’s state.
To clarify, the component distributions P (k | z=j) = Ωij qj depend on the variational
posteriors qj of other agents, making the generative model itself a function of variational
quantities. This construction can be understood via two complementary points. First, as
an augmented model whereby the full system has a joint generative model over all agents’
latent states, and the pairwise KL terms emerge as variational couplings in the mean-field
decomposition of the joint ELBO. Second, as a consensus energy whereby the alignment
free energy P
j βij DKL(qi∥Ωij qj ) can be viewed as an engineered energy functional that
penalizes belief disagreement after gauge transport, motivated by the KL term of a fixed
generative model. In either framing, the resulting update equations are the same, and the
key property is that at a fixed point of the joint optimization (over all {qi} simultaneously),
the pairwise objectives are mutually consistent.
8
Attention and Transformers as Free Energy Minimization
3.4.2 Variational Posterior
Next, we approximate the true posterior P (k, z | data) using mean-field factorization:
Q(k, z) = qi(k) · β(z), (16)
where β(z = j) ≡ βij ≥ 0 are variational attention weights satisfying P
j βij = 1.
3.4.3 Alignment Free Energy
The free energy for the alignment component is the KL divergence between the variational
posterior and the generative model:
Falign = DKL[Q(k, z)∥P (k, z)] = X
j
βij
Z
dk qi(k) log qi(k) βij
P (k | z=j) πj
. (17)
Decomposing the logarithm yields:
Falign = X
j
βij





Z
dk qi(k) log qi(k)
P (k | z=j)
| {z }
= DKL[qi ∥ Ωij qj ]
+ log βij − log πj




 . (18)
The identification of the integral with DKL[qi∥Ωij qj ] follows directly from the definition
of KL divergence and the generative model P (k | z=j) = (Ωij qj )(k), with no approximation.
Therefore, we may define an alignment energy Eij ≡ DKL[qi∥Ωij qj ] such that:
Falign = X
j
βij
Eij + log βij − log πj
. (19)
This has the usual ”energy minus entropy” form Falign = ⟨E⟩β − H(β) + const, where
H(β) = − P
j βij log βij is the entropy of the attention distribution.
3.4.4 Minimization and the Softmax Solution
Next, we minimize Falign with respect to βij subject to the constraint P
j βij = 1 using
Lagrange multipliers:
L = X
j
βij (Eij + log βij − log πj ) − λ

X
j
βij − 1

 . (20)
Requiring ∂L/∂βik = 0 produces:
Eik + log βik + 1 − log πk − λ = 0. (21)
Solving and normalizing we find the softmax form (importantly) weighted by attention
priors πk:
βik = πk exp(−Eik)
P
m πm exp(−Eim) . (22)
9
Dennis
For example, a uniform attention prior πk = 1/N , the prior factors completely cancel
yielding:
βik = softmaxk(−Eik) = softmaxk
 − DKL[qi∥Ωikqk]. (23)
Without loss of generality we may include a temperature parameter τ > 0 by rescaling
the alignment energy Eik → Eik/τ . Additionally, as discussed in the appendix, the forward
KL divergence is the unique f -divergence that preserves exponential-family closure under
this linear coupling structure and yields a consistent dual interpretation for the attention
weights (see Appendix C).
3.4.5 Non-Uniform Attention Priors: Masking and Positional Biases
The general solution (22) with non-uniform priors πk provides a unified derivation of several
architectural features historically introduced empirically into transformers.
Causal Masking. In auto-regressive language modeling, token i should not attend to
future tokens j > i. Within our generative model, this constraint is a prior restriction on
the availability of a source. Formally, setting the attention prior to
π(causal)
j =
( 1
i if j ≤ i,
0 if j > i, (24)
and substituting into (22) produces
βij = 1[j ≤ i] · exp(−Eij /τ )
P
k≤i exp(−Eik/τ ) , (25)
which is precisely the masked attention used in GPT-family models (Radford et al.,
2019).
Sliding Window Attention. Similarly, choosing the prior as a local window of size w
around position i,
π(window)
j ∝ 1[|i − j| ≤ w], (26)
produces the sliding window attention used in architectures such as Longformer (Belt-
agy et al., 2020) and Mistral (Jiang et al., 2023). In our gauge-theoretic framework, this
encodes the prior belief that an agent’s state was generated by a source within its ”local
neighborhood”, a natural assumption for data with locality structure.
ALiBi: Linear Positional Bias. A distance-decaying attention prior of the form
πj ∝ exp(−m|i − j|), (27)
where m > 0 is a slope parameter, contributes an additive bias to the logits. Substituting
into (22) we have:
βij = exp(−Eij /τ − m|i − j|)
P
k exp(−Eik/τ − m|i − k|) = softmaxj

− Eij
τ − m|i − j|

. (28)
10
Attention and Transformers as Free Energy Minimization
This is exactly the ”Attention with Linear Biases” (ALiBi) mechanism of Press et al.
(2022). The slope m is typically set to a head-dependent constant (e.g., m = 2−8/H · 2−h
for head h of H total heads). In the gauge-theoretic framework, ALiBi arises from a prior
belief that nearby sources are exponentially more likely to have generated the agent’s state.
Relative Position Biases. More generally, we can make the attention prior a learnable
function of relative position,
πj ∝ exp(bi−j ), (29)
where bδ ∈ R is a (learnable) bias for relative offset δ = i − j, leading to an additive bias
bi−j in the logits as:
βij = softmaxj

− Eij
τ + bi−j

. (30)
This recovers relative position biases used in T5 (Raffel et al., 2020) and other architec-
tures that parameterize attention with learnable position-dependent offsets. In a straight-
forward manner, each positioning architecture emerges as a special case of (22) operating
via the same mathematical object. Although the choice of attention priors may be ad-
hoc, their fundamental source and structure is determined from first principles within our
framework.
3.5 Full Variational Free Energy
By combining the single-agent free energy (14) with the alignment terms derived above and
summing over all agents, we obtain the global gauge-covariant variational free energy at
c∗ ∈ C:
F[{qi}] = X
i
DKL(qi ∥ pi) + X
i,j
β∗
ij DKL(qi ∥ Ωij qj ) − Eq[log p(o |{ki})]. (31)
We note that this is the partially minimized free energy: the attention weights β∗
ij have
already been set to their optimal softmax values (22) by minimizing the full alignment free
energy Falign (19) with respect to β. The full (un-minimized) alignment energy includes an
attention entropy/prior term τ P
i,j βij log(βij /πj ) (cf. Eq. 19), which is essential for it is
this term that makes the β-optimization yield the softmax rather than an arg-min. When
βij is set to its optimal value β∗
ij = softmaxj (−Eij /τ +log πj ), the entropy term evaluates to
−τ P
i log Zi (where Zi = P
j πj exp(−Eij /τ ) is the partition function), which is absorbed
into F[{qi}] as a function of beliefs alone. This absorbed log-partition function does not
appear as an explicit loss term in the implementation, but it contributes to the gradients
∂F/∂μi through the chain rule via ∂β∗/∂μi representing a fundamental non-linearity in the
theory as we dkiscuss further below.
The envelope theorem implies consistency since:
∂F/∂μi β=β∗ = d
dμi
F[{qi}, β∗(q)]
since the cross-term
11
Dennis
X
j
(∂F/∂βij )(∂β∗
ij /∂μi)
vanishes at the optimal attention. The gradient expressions in Section 4.3, which include
∂βij /∂μi terms, are therefore correct as written under this partially minimized convention
(similarly for Σ and ϕ) .
3.6 Interpretation
Each term in the variational free energy (31) encodes a distinct aspect of multi-agent infer-
ence:
(1) Belief Prior: DKL(qi∥pi)
The belief prior measures the deviation of agent i’s current belief qi(ki) from its local
prior pi(ki | mi). This term regularizes beliefs towards its expected states. Without obser-
vations and inter-agent couplings, minimizing F would produce q∗
i = pi thereby recovering
pure prior-based prediction.
(2) Belief Alignment: βij DKL(qi∥Ωij qj )
As discussed previously, belief alignment represents the discrepancy between agent i’s
belief and agent j’s belief after gauge transport into agent i’s local frame (i.e. agent i’s
interpretation of agent j’s belief).
This term enforces an epistemic consensus whereby agents with high βij are driven to
agree on their beliefs about the current state, modulo gauge transformations.
(3) Observation Likelihood: −Eq[log p(o | {ki})]
This term is the expected negative log-likelihood of observations/data o given latent
states {ki}, averaged over the recognition distributions {qi}. Without this term, the system
is a pure vacuum theory where agents converge to their coupled priors without external
input. Observations break this vacuum symmetry, forcing agents to specialize based on
their sensory evidence (see Results).
3.7 State-Dependent Prior Precision
3.7.1 Precision as a Variational Parameter
In the variational free energy (31), the self-coupling term DKL(qi∥pi) carries unit coefficient.
This implicitly assumes that the prior precision is apriori known and fixed such that each
agent trusts its prior with uniform strength regardless of context.
The hierarchical structure of the generative model (discussed in the appendix) implies
that the effective prior precision manifestly depends on the model uncertainty. Therefore,
we may choose to promote the self-coupling weight from a fixed constant to a per-agent
variational parameter αi > 0 (subject perhaps to a regularizer that prevents degenerate
solutions).
This augmented free energy for agent i is then:
Fi = αi DKL(qi∥pi) + R(αi) + X
j
βij DKL(qi∥Ωij qj ) − Eqi [log p(oi | ki)], (32)
where R(αi) is a precision regularizer. A natural choice is a log-barrier form such as:
12
Attention and Transformers as Free Energy Minimization
R(αi) = b0 αi − c0 log αi, b0 > 0, c0 > 0. (33)
The first term penalizes arbitrarily large precision (preventing αi → ∞), the second is
a log-barrier that enforces αi > 0 and penalizes vanishing precision (preventing αi → 0,
which would trivially eliminate the self-coupling cost).
3.7.2 Optimal Precision
Optimizing Fi with respect to αi (holding beliefs qi fixed):
∂Fi
∂αi
= DKL(qi∥pi) + b0 − c0
αi
= 0, (34)
which gives the optimal state-dependent precision:
α∗
i = c0
b0 + DKL(qi∥pi) . (35)
For example, when an agent is close to its prior (DKL(qi∥pi) ≈ 0), αi ≈ c0/b0 takes a
large value determined by the hyper-parameters such that the agent trusts its prior and
resists displacement by the alignment term. Conversely, when an agent is far from its prior
(large DKL(qi∥pi)), αi → 0, the prior’s influence is reduced and the alignment and obser-
vation terms dominate. The agent becomes less influenced by its prior allowing contextual
information to influence its belief.
b0 controls the sensitivity such that large b0 yield approximately constant αi, while small
b0 yield sharply state dependent αi. The ratio c0/b0 sets the maximum precision (the value
when beliefs match the prior exactly). Both hyper-parameters may be learned via empirical
Bayes by introducing two scalar parameters.
In our experiments we study both R(αi)̸ = 0 and R(αi) = 0 with αi = α = 1.
3.8 Belief Dynamics
The global free energy (31) is minimized by gradient descent on the beliefs {qi}:
∂qi
∂t = −ηq
δF
δqi
(36)
The most general formulation includes additional model-channel terms (si, ri, γij ) rep-
resenting a slow meta-learning timescale (Appendix F.1). In the present work we focus
exclusively on the belief (fast) subsystem, studying how beliefs {qi} evolve under align-
ment. We posit that standard transformers operate entirely in this regime: they perform
inference (qi updates) with frozen priors during a forward pass. As discussed in the ap-
pendix, the framework naturally supports hierarchical emergence of meta-agents through
coarse-graining when groups of agents reach belief consensus (Appendix A.2.1).
Note that in the full theory with general, possibly curved, base manifolds we would
necessarily integrate over agent supports and overlaps in order to obtain the complete vari-
ational free energy (which we needn’t consider for our present transformer derivations).
This then introduces extensive and rich geometry including Riemann curvature, gauge cur-
vature, Fischer curvature, induced gauge connections, belief/model transport between and
13
Dennis
within agents and fibers, non-trivial vertical and horizontal holonomies, and much, much
more.
3.9 Symmetry Breaking
In the absence of observations (p(o|·) = const), the free energy (31) possesses a reparame-
terization degeneracy under the gauge group G (which is GL(K) in general, or a compact
subgroup such as SO(K)). The vacuum state corresponds to perfect alignment: qi = Ωij qj
and pi = ˜Ωij pj for all agent pairs (i, j), meaning that all agents maintain gauge-equivalent
beliefs that differ only by frame transformations. In this regime, the dynamics drive the
system toward a degenerate manifold of ground states parameterized by the gauge orbit.
Agents synchronize over gradient descent towards ∥μi(c)∥ = μ∗, but the absolute orienta-
tions remain arbitrary.
A terminological note is warranted. In physics, local gauge symmetries are more pre-
cisely understood as redundancies of description rather than physical symmetries, and there
are well-known subtleties in applying the language of spontaneous symmetry breaking to
gauge redundancies (notably, Elitzur’s theorem forbids spontaneous breaking of local gauge
symmetries in lattice gauge theory (Weinberg, 1996)). Our gauge group acts as a reparam-
eterization symmetry of the KL divergence (Theorem 1), and the degeneracy of the vacuum
manifold is a consequence of this reparameterization freedom rather than of a global physical
symmetry.
Observations destroy this degeneracy by coupling agents to external data through the
likelihood terms. Each agent’s data stream oi acts as an external field that selects a preferred
orientation in its fiber. This is explicit symmetry breaking: the observation/likelihood term
−Eq[log p(o|k)] is manifestly not gauge-invariant, and it is this term that forces agents into
distinct configurations. The system transitions from the symmetric vacuum to a phase
where agents develop distinct specializations such that ∥μi(c)∦ = ∥μj (c)∥ with the diver-
sity driven by observations/data. The frame transformations Ωij then encode how these
specialized representations relate geometrically. This observation-induced specialization is
what enables non-trivial multi-agent coordination in that agents must actively maintain
geometric relationships between their specialized frames rather than simply coexisting in a
gauge-symmetric configuration. Whether the vacuum manifold supports approximate zero
modes of the Hessian (corresponding to directions of near-zero curvature along the gauge
orbit) is a question we leave to future work.
3.10 Summary
The gauge-covariant variational free energy (31) has been derived from two principles: (i)
the standard free energy principle for single agents (prior regularization + observation like-
lihood), and (ii) the mixture-of-sources model for inter-agent communication (KL-based
alignment with softmax attention weights). Both belief alignment and model alignment
arise from variational inference over gauge-transported sources. The forward KL diver-
gence DKL(qi∥Ωij qj ) emerges as the alignment energy from the mixture generative model,
and the softmax form of attention weights βij = softmax(−Eij /τ ) follow from Lagrange op-
timization of the alignment free energy. Non-uniform attention priors πk provide multiple
theoretically grounded mechanisms for positional relationships. Furthermore, a context-
14
Attention and Transformers as Free Energy Minimization
dependent self-coupling weight αi can be derived as a state-dependent precision through
variational optimization with a log-barrier regularizer (Section 3.7), with the fixed-α case
recovered as a limit.
Multi-agent communication within this gauge-covariant formulation allows the FEP to
be satisfied. Importantly, there is no necessitation of any informational substrate (such
as neural networks, transmission lines, etc) within our formalism. Next, we show that
this formulation connects attention, transformers, and architecture to variational inference
without neural networks, MLPs, and other standard ad-hoc architectures.
4 Reduction to Transformer Attention
In this section we demonstrate that standard transformer self-attention emerges as a set of
limiting cases of our gauge-theoretic framework. This establishes neural network training as
a limit of a first-principles information-geometric minimization of a variational free energy
over a bundle geometry.
4.1 Setup: Agents as Gaussian Beliefs in Local Frames
In our general formulation, each agent i (token) maintains a local MVG belief:
qi = N (μi, Σi), (37)
where μi ∈ Rdk is the agent’s mean representation in its local gauge frame, and Σi ∈
Rdk ×dk is its belief covariance (symmetric positive definite). Note that we are embedding
in K = dk dimensions in order to make contact with standard notation and, for now, the
gauge group will be the fundamental GL(K = dk).
Communication between agents i and j is mediated by a gauge frame transport operator
Ωij ∈ GL(dk), (38)
which maps representations from agent j’s local frame into agent i’s frame.
The message (update) received by agent i is then
mi = X
j
βij Ωij μj , (39)
which is the gauge-theoretic analog of the attention aggregation P
j αij Vj in standard
transformers.
4.2 Connection to Standard Transformers
The gauge-theoretic framework is inherently more general than standard transformer at-
tention. To establish the connections formally, we show that successive limits reduce the
full gauge-covariant attention to the familiar dot-product form. Our experiments, however,
(Section 5.3), implement the most general model.
15
Dennis
Deterministic Beliefs via Scaled Limit. We consider a single-point base space C with
a finite set of agents {1, . . . , N } (tokens) with no particular ordering.
The deterministic limit requires nuance. Naively setting Σi → 0 yields Dirac delta
beliefs qi(ki) → δ(ki − μi), but the KL divergence between distinct Diracs is infinite:
DKL(δ(k − μi)∥δ(k − μj )) = +∞ for μi̸ = μj . (40)
This divergence reflects the absolute continuity requirement in the definition of KL:
DKL(P ∥Q) is finite only when P is absolutely continuous with respect to Q (i.e., P (A) > 0
implies Q(A) > 0 for all measurable A). Distinct Dirac measures have disjoint support and
therefore violate this condition.
We may remedy this by taking a joint scaling limit where the belief variance σ2 (discussed
below) remains finite but is absorbed into learned parameters.
Next, we assume isotropic covariances Σi = Σj = σ2I with σ2 > 0. This implies that the
log-determinant terms cancel (log |Σj |
|Σi| = 0) and the trace term simplifies to Tr(Σ−1
j Σi) =
Tr(I) = dk.
With gauge transport Ωij ∈ GL(dk), the transported Gaussian Ωij qj has mean Ωij μj
and covariance Ωij Σj Ω⊤
ij = σ2(Ωij Ω⊤
ij ). Therefore the full KL divergence is:
DKL(qi∥Ωij qj ) = 1
2
h
log det(Ωij Ω⊤
ij )+Tr(Ωij Ω⊤
ij )−1−dk
i
+ 1
2σ2 (μi−Ωij μj )⊤(Ωij Ω⊤
ij )−1(μi−Ωij μj ).
(41)
The Mahalanobis quadratic form can be simplified using the identity (Ωij Ω⊤
ij )−1 =
Ω−⊤
ij Ω−1
ij :
(μi − Ωij μj )⊤(Ωij Ω⊤
ij )−1(μi − Ωij μj ) = ∥Ω−1
ij (μi − Ωij μj )∥2 = ∥Ω−1
ij μi − μj ∥2. (42)
This identity has the interpretation that rather than transporting the key belief into the
query’s frame via Ωij and measuring with the induced metric, the KL divergence equiva-
lently transports the query into the key’s frame via Ω−1
ij and measures the Euclidean distance
there. Either perspective is appropriate. The Mahalanobis metric automatically accounts
for the distortion introduced by gauge transport.
The structural terms log det(Ωij Ω⊤
ij ), Tr((Ωij Ω⊤
ij )−1), and dk depend on the gauge trans-
port but not on the belief means μi, μj directly. Under Limit 2 (Ωij = Ω), they become
constant across all agent pairs and cancel under softmax thereby producing:
DKL(qi∥Ωqj ) ∼ 1
2σ2 ∥Ω−1μi − μj ∥2 + (terms absorbed by softmax). (43)
Next, the rescaled KL divergence limit leads to:
eDKL(qi∥Ωij qj ) := σ2 · DKL(qi∥Ωij qj ) σ→0
−−−−→ 1
2 ∥Ω−1
ij μi − μj ∥2, (44)
which remains finite and equals half the squared distance between means as measured
in the key’s reference frame. This rescaled quantity therefore has a well-defined limit as
σ → 0 and represents the natural information-geometric measure of dis-similarity within the
16
Attention and Transformers as Free Energy Minimization
deterministic regime. This should be understood as the regime where σ is small compared
to characteristic scales of μi, such that σ−2 will be absorbed into the learned projections
(see below).
Flat Bundle Limit (Ωij = Ω). We next assume a single global frame shared by all
agents; i.e. a flat connection with no position-dependent frame misalignment:
Ωij = Ω ∈ Rdk ×dk for all i, j, (45)
with Ω being a fixed linear map.
4.2.1 Derivation of Dot-Product Attention
Expanding the KL Divergence. After our limits the KL divergence reduces to (drop-
ping constant terms that cancel under softmax):
sij = DKL(qi∥Ωqj ) = 1
2σ2 ∥Ω−1μi − μj ∥2 + C, (46)
where C collects the remaining terms (log-determinant, trace, dk) that are independent
of the belief means and ultimately cancel under softmax. Expanding the squared norm we
have:
sij = 1
2σ2 ∥Ω−1μi∥2 + 1
2σ2 ∥μj ∥2 − 1
σ2 μ⊤
i Ω−⊤μj + C. (47)
Softmax Cancellations. For a fixed query agent i, the softmax over keys j is:
βij = exp(−sij /τ )
P
k exp(−sik/τ ) . (48)
The i-terms 1
2σ2 ∥Ω−1μi∥2 + C are independent of j and therefore cancel between numer-
ator and denominator of the softmax yielding:
βij ∝ exp
 1
τ
 1
σ2 μ⊤
i Ω−⊤μj − 1
2σ2 ∥μj ∥2

. (49)
Thus, the effective logit is:
˜sij = 1
σ2 μ⊤
i Ω−⊤μj − 1
2σ2 ∥μj ∥2. (50)
Key Bias Cancellation. The second term − 1
2σ2 ∥μj ∥2 however, depends on key j and is
a key-dependent bias on the raw embedding norms. This bias involves ∥μj ∥2 rather than
∥Ωμj ∥2 since the Mahalanobis metric (42) already accounts for the gauge transport’s effect
on norms via the (ΩΩ⊤)−1 factor in the distance measure. This residual bias arises purely
from the embedding space itself. Two well-trodden mechanisms ensure its approximate
cancellation:
(1) High-Dimensional Concentration: For embeddings in Rdk with approximately
independent components (e.g., μ(i)
j ∼ N (0, σ2
0 )), the concentration of measure implies:
17
Dennis
∥μj ∥2 =
dkX
i=1
(μ(i)
j )2 = dkσ2
0 ± O(σ2
0
pdk), (51)
with relative fluctuations O(1/√dk) → 0 as dk increases. Note: this σ0 has nothing to
do with σ
(2) Layer Normalization: Modern transformer architectures enforce constant key
norms explicitly via layer normalization, which normalizes ∥μj ∥ to be constant across tokens.
This makes the key bias:
− 1
2σ2 ∥μj ∥2 ≈ (constant in j), (52)
which cancels under softmax.
However, on the gauge theory side, the algebraic origin of this cancellation can be
understood partly through the Lie algebraic decomposition of the gauge group. The Lie
algebra gl(K) splits as a direct sum:
gl(K) = sl(K) ⊕ R, (53)
where sl(K) consists of traceless matrices (volume-preserving transformations) and R
is the trace component (overall scaling). The key-norm bias ∥μj ∥2 depends on the magni-
tude of the embedding vector. Layer normalization, which centers and rescales μj to have
constant norm, effectively quotients out the positive scaling degree of freedom from the
embedding space, projecting representations onto a sub-manifold of constant norm. This
eliminates the source of key-norm bias at its algebraic origin. Combined with the Maha-
lanobis metric (which handles the gauge transport’s effect on norms through (ΩΩ⊤)−1), this
ensures that attention depends on belief content rather than representation magnitude.
Factorization into Query-Key Dot Product. The leading term in ˜sij is:
1
σ2 μ⊤
i Ω−⊤μj . (54)
We identify the product of WQ, WK with the gauge transport:
WQW ⊤
K = 1
σ2 Ω−⊤ ∈ GL(dk). (55)
Such a factorization always exists for any invertible Ω since for Ω−⊤ ∈ GL(dk), any
matrix M ∈ GL(dk) can be written as M = AB⊤ for suitable A, B (e.g., via SVD: M =
U ΛV ⊤ = (U Λ1/2)(V Λ1/2)⊤).
Importantly this identification concerns the bilinear form M = WQW ⊤
K that determines
the attention logits rather than the individual matrices WQ and WK separately. In standard
transformer implementations, the per-head projection matrices W a
Q, W a
K ∈ Rdmodel×dhead are
rectangular (with dmodel ≫ dhead), and neither is individually invertible in the dmodel-
dimensional space. The gauge-theoretic identification is within the dhead-dimensional pro-
jected subspace where the effective bilinear form M a = (W a
Q)⊤W a
K ∈ Rdhead×dhead acts.
The standard rectangular shape of W a
Q arises from a further restriction: decomposing Ω
into a block-diagonal subgroup GL(dhead)H ⊂ GL(dk), where each head a operates within
18
Attention and Transformers as Free Energy Minimization
an independent dhead-dimensional subspace. The rectangular shape encodes both the or-
thogonal projection onto block a and the GL(dhead) transport within that block. We develop
and discuss this identification in detail in Section 4.4. For now we return to the attention
derivation.
Next, rather than completely taking σ → 0 we recognize that σ−2 and Ω−⊤ always
appear together in the combination σ−2Ω−⊤. The learned matrices WQ, WK can then be
considered to parametrize this combined quantity directly, rendering σ an implicit (finite)
scale factor absorbed into the learned weights. Therefore, the full limit need not be taken.
Next, we define query and key vectors as:
Qi ≡ μ⊤
i WQ ∈ R1×dk , Kj ≡ μ⊤
j WK ∈ R1×dk . (56)
Then:
QiK⊤
j = μ⊤
i WQW ⊤
K μj = 1
σ2 μ⊤
i Ω−⊤μj . (57)
Thus the logit term reduces to a standard dot product as:
˜sij = QiK⊤
j + O(1), (58)
where O(1) represents the approximately constant key bias (controlled by layer normal-
ization).
Temperature Scaling. Since σ−2 is absorbed into the learned projections WQ, WK ,
the effective logit ˜sij = QiK⊤
j carries no residual dependence on σ. The softmax βij =
softmaxj (˜sij /τ ) must therefore set τ to account only for the dimensional scaling of the dot
product.
In high-dimensional spaces, dot products QiK⊤
j have standard deviation O(√dk) (for
unit-variance components), while key-bias fluctuations are O(1) after layer normalization.
Normalizing pre-softmax logits to O(1) variance therefore requires:
τ = pdk. (59)
Combining all simplifications we have:
βij = softmaxj
QiK⊤
j
√dk
!
, (60)
exactly recovering the standard transformer attention weighting rule.
4.2.2 Value Aggregation
The message aggregation rule is:
mi = X
j
βij Ωij μj . (61)
19
Dennis
GL(K) aggregation and the role of Ω−⊤. The mixture-of-sources generative model
(Section 3.4) specifies that the component distribution for source j in agent i’s frame has
mean Ωij μj . The posterior mean under mixture responsibilities βij is therefore:
mi = X
j
βij Ωij μj , (62)
which is the standard mixture-of-Gaussians expectation. This is the aggregation con-
sistent with the generative model..
A separate geometric object arises in the attention weight computation, where the Ma-
halanobis identity (42) introduces Ω−⊤ as the natural transform for the dual (covector)
structure of the inner product. Specifically, Ω−⊤ appears in the logit μ⊤
i Ω−⊤μj because the
KL divergence measures distance in the transported covariance metric (ΩΩ⊤)−1, and Ω−⊤
is the correct transform for covectors (dual vectors) under a linear map. To be clear, this
dual transform belongs to the attention scoring geometry, not to the value aggregation.
Next, we define the value projection:
Vj ≡ W ⊤
V μj , WV ∈ Rdk ×dk , (63)
where we absorb the gauge transport Ω into the learned matrix WV .
Then:
mi = X
j
βij Vj , (64)
identical to the standard transformer attention update zi = P
j αij Vj .
4.2.3 Complete Attention Formula
Under our assumptions, framework, and limits, our gauge-theoretic message communication:
βij ∝ exp

− 1
τ DKL(qi∥Ωij qj )

, mi = X
j
βij Ωij μj , (65)
reduces to:
Attention(Q, K, V ) = softmax
 QK⊤
√dk

V . (66)
This is the canonical dot-product attention mechanism in transformers.
The limits are deliberately aggressive. They collapse the statistical manifold into a basic
Euclidean space and absorbing the gauge parameters into the learned matrices.
Each limit discards a specific geometric degree of freedom and each can be independently
relaxed to produce a generalization of standard attention. The gauge-theoretic framework
provides a taxonomy of design choices whereby standard attention corresponds to the three-
limit simplification, while much richer architectures (such as non-isotropic beliefs, learned
gauge transport, position-dependent connections, etc) occupy intermediate points which
may have unique and interesting applications.
20
Attention and Transformers as Free Energy Minimization
Our GL(K) experiments (Section 5.3) utilize the full generality (at least for a 0 di-
mensional manifold) and demonstrate that meaningful language modeling is possible using
the complete gauge-theoretic variational free energy alone without incorporating standard
neural network architectures.
4.3 Training as Variational Free Energy Minimization
Having established that attention emerges from the variational free energy, we now show
that gradient descent on this free energy recovers standard transformer training dynamics,
providing a variational interpretation of each component.
An important distinction must be maintained between two operations that the free
energy framework treats uniformly but that standard transformers implement differently.
In variational inference, minimizing F with respect to the variational density q constitutes
inference (the E-step), while updating generative model parameters constitutes learning
(the M-step). In a standard transformer, the forward pass computes per-token hidden states
via deterministic layer maps (analogous to amortized inference whereby a single learned
function replaces iterative optimization of q), while gradient descent on weights across data
batches constitutes learning. The gradient expressions below describe the structure of both
operations since they share the same functional form. However, what differs is in which
variables are updated and on what timescale. The forward pass corresponds to a single
step of belief inference with frozen parameters, and weight training corresponds to slow
parameter updates that reduce the expected free energy across data.
4.3.1 Observation Likelihood as Loss Function.
In our gauge theory, observations by agents act as a source term that breaks the vacuum
degeneracy. The vacuum free energy (without observations), written in partially minimized
form with βij at its optimal softmax value (cf. Eq. 31), is:
Fvacuum[{qi}] = X
i
αiDKL(qi∥pi) + X
i,j
βij DKL(qi∥Ωij qj ). (67)
Introducing per-agent observations oi adds the likelihood term:
F[{qi}] = Fvacuum[{qi}] − X
i
Eqi [log p(oi | ki)]. (68)
In the deterministic limit qi → δ(ki − μi), this becomes:
−Eqi [log p(oi | ki)] → − log p(oi | μi), (69)
the standard negative log-likelihood.
For Gaussian observations p(o | μ) = N (o | μ, Σ):
− log p(o | μ) = 1
2 (o − μ)⊤Σ−1(o − μ) + const. (70)
=⇒ Lobs = 1
2 ∥o − μ∥2
Σ (Mean-squared error). (71)
21
Dennis
For categorical observations p(o | μ) = Categorical(softmax(μ)):
− log p(o | μ) = − X
k
ok log(softmax(μ)k) (Cross-entropy loss). (72)
These are standard loss functions in machine learning.
Deterministic Free Energy. In the deterministic limit with flat bundle (Ωij = Ω), each
belief becomes a point qi → δ(ki − μi) and each prior is isotropic: pi = N (μprior, σ2
p I). The
self-coupling term from the adaptive free energy (32) becomes:
αi DKL(δ(μi) ∥ N (μprior, σ2
p I)) = αi
2σ2
p
∥μi − μprior∥2, (73)
where the state-dependent precision (35) simplifies in the deterministic limit (with
DKL(δ(μi)∥N (μprior, σ2
p I)) = 1
2σ2
p
∥μi − μprior∥2 and absorbing constants) to:
αi = c0
b0 + 1
2σ2
p
∥μi − μprior∥2
. (74)
Defining the per-token regularization strength λp,i ≡ αi/σ2
p , the full deterministic free
energy is:
F[{μi}] = X
i
λp,i
2 ∥μi − μprior∥2
| {z }
Adaptive prior regularization
+ X
i,j
βij
2σ2 ∥Ω−1μi − μj ∥2
| {z }
Alignment coupling
− X
i
log p(oi | μi)
| {z }
Observation likelihood
. (75)
As previously discussed this form constitutes an adaptive weight decay whereby tokens
whose beliefs have drifted far from the prior (∥δi∥ large) experience weaker regularization
(λp,i small, gate open), while tokens near the prior are strongly anchored (λp,i large, gate
closed).
In the limit where αi → α = const and λp,i → λp ≡ α/σ2
p we recover uniform L2
regularization:
F[{μi}] = X
i
λp
2 ∥μi − μprior∥2
| {z }
L2 regularization (weight decay)
+ X
i,j
βij
2σ2 ∥Ω−1μi − μj ∥2
| {z }
Alignment coupling
− X
i
log p(oi | μi)
| {z }
Observation likelihood
. (76)
This identifies the standard weight decay coefficient λp as the product of two quantities
that the general theory treats independently: the prior precision α and the inverse prior
variance 1/σ2
p . In the full theory, both are generally state-dependent.
22
Attention and Transformers as Free Energy Minimization
4.3.2 Gradient Descent Dynamics.
In the limit λp,i = λp the gradient of the free energy with respect to μi is:
∂F
∂μi
= λp (μi − μprior)
+ X
j
 βij
σ2 (ΩΩ⊤)−1 (μi − Ωμj ) + 1
2σ2
∂βij
∂μi
Ω−1μi − μj
2

− ∂ log p(oi | μi)
∂μi
.
(77)
The first term, λp(μi − μprior), as discussed, matches weight decay or L2 regularization.
The second term, P
j
βij
σ2 (ΩΩ⊤)−1(μi −Ωμj ), is attention-weighted message aggregation that
pulls μi toward gauge-transported neighbors. The third term, P
j
∂βij
∂μi ∥Ω−1μi − μj ∥2/(2σ2),
represents gradient flow through the attention weights themselves - a manifestly non-linear
”activation”. The fourth term, − ∂ log p(oi|μi)
∂μi , is the gradient of the loss function (i.e. the
supervised learning signal).
Since beliefs qi live on the Gaussian statistical manifold, the geometrically correct dy-
namics uses the natural gradient. The Euclidean gradient must be pre-conditioned by the
inverse Fisher information metric which, for MVGs is G−1
μ = Σi (see Appendix D.2 for
derivation). The natural gradient flow is:
dμi
dt = −η Σi
∂F
∂μi
. (78)
In the isotropic limit Σi = σ2I, the Fisher preconditioning absorbs the σ−2 factors in
the alignment terms. Defining the effective learning rate ˜η ≡ ησ2 and discretizing yields:
μ(t+1)
i = μ(t)
i − ˜η
"
λp (μi − μprior) + X
j

βij (ΩΩ⊤)−1 (μi − Ωμj ) + 1
2
∂βij
∂μi
Ω−1μi − μj
2

− ∂ log p(oi | μi)
∂μi
#
.
(79)
The natural gradient recovers the standard transformer update rule. The alignment
term P
j βij Ωμj maps directly to the value aggregation P
j αij Vj without spurious σ−2
factors. This is consistent with the E-step algorithm we implement, which applies Fisher
preconditioning ˜∇μ ← Σi · ∇μF.
General (non-isotropic) mean gradient. The preceding expressions assumed isotropic
covariances Σi = σ2I and a shared gauge connection Ω. We now state the fully general
gradient, valid for arbitrary covariances Σi, edge-dependent gauge connections Ωij , general
prior precision Σ−1
p,i , and state-dependent prior coupling αi (Section 3.7).
Differentiating the full free energy (31) with respect to μi and applying the product rule
through both the attention weights βij (μi) and the adaptive prior coupling αi(μi) gives:
23
Dennis
∂F
∂μi
= αi Σ−1
p,i (μi − μprior) + ∂αi
∂μi
DKL(qi ∥ pi)
+ X
j

βij

Ωij Σj Ω⊤
ij
−1
(μi − Ωij μj ) + ∂βij
∂μi
DKL(qi ∥ Ωij qj )

− ∂ log p(oi | μi)
∂μi
.
(80)
The structure parallels the isotropic case term-by-term, with four generalizations: (i) the
prior precision λpI is replaced by the full matrix αi Σ−1
p,i , where αi is the state-dependent
prior coupling; (ii) the term (∂αi/∂μi) DKL(qi ∥ pi) represents gradient flow through the
adaptive gate — agents far from their prior (large ∥μi − μprior∥) have ∂αi/∂μi < 0, which
reduces the effective regularization and allows beliefs to drift further, reinforcing the gating
behavior described in Section 3.7; (iii) the isotropic transported precision σ−2(ΩΩ⊤)−1 is
replaced by the edge-dependent transported precision (Ωij Σj Ω⊤
ij )−1, which accounts for
both the neighbor’s covariance structure and the geometry of the gauge transport along
each edge; and (iv) the quadratic mean-mismatch 1
2σ2 ∥Ω−1μi − μj ∥2 is replaced by the
full KL divergence DKL(qi ∥ Ωij qj ), which additionally couples to the covariance mismatch
between Σi and the transported covariance Ωij Σj Ω⊤
ij . The full KL cannot be reduced to
a purely mean-dependent expression because the covariance contributions (trace and log-
determinant terms) are j-dependent through Ωij and Σj , and therefore do not cancel under
the constraint P
j ∂βij /∂μi = 0.
When αi = α = const, the gate gradient vanishes (∂αi/∂μi = 0) and the prior term
reduces to α Σ−1
p,i (μi − μprior), recovering the uniform regularization case.
The corresponding natural gradient flow is:
dμi
dt = −η Σi
∂F
∂μi
= −η
"
ΣiΣ−1
p,i (μi − μprior) + X
j

βij Σi

Ωij Σj Ω⊤
ij
−1
(μi − Ωij μj )
+ ∂βij
∂μi
Σi DKL(qi ∥ Ωij qj )

− Σi
∂ log p(oi | μi)
∂μi
#
.
(81)
When the covariance alignment condition Σi ≈ Ωij Σj Ω⊤
ij holds (as enforced by the
covariance gradient dynamics below), the Fisher-preconditioned alignment term simplifies to
βij (μi−Ωij μj ), recovering attention-weighted message passing with geometric transport (the
non-isotropic generalization of the value aggregation P
j αij Vj ). This covariance alignment
condition is precisely the fixed point of the Σ-gradient (Eq. 83), demonstrating internal
consistency whereby the mean and covariance dynamics cooperate to produce clean message
passing at equilibrium.
Covariance gradient. Similarly, the complete theory also includes gradient dynamics for
the covariance Σi and gauge frame ϕi. The covariance gradient reveals a novel uncertainty
channel absent in standard transformers. Differentiating the free energy with respect to Σi
24
Attention and Transformers as Free Energy Minimization
and applying the product rule through both the attention weights and the adaptive prior
coupling αi(Σi) (neglecting the observation term, which contributes an O(Σ−1
i ) correction
in the high-precision regime; see Appendix B for the full derivation) yields:
∂Fi
∂Σi
= 1
2

−(1 + αi) Σ−1
i + αi Σ−1
p,i + X
j
βij (Ωij Σj Ω⊤
ij )−1


+ ∂αi
∂Σi
DKL(qi ∥ pi) + X
j
∂βij
∂Σi
DKL(qi ∥ Ωij qj ).
(82)
This parallels the mean gradient term-by-term: αi Σ−1
p,i is the prior precision scaled by
the adaptive coupling (analogous to weight decay for μ); the (∂αi/∂Σi) DKL(qi ∥ pi) term is
gradient flow through the adaptive gate (the covariance analogue of the (∂αi/∂μi) DKL term
in Eq. 80); P
j βij (Ωij Σj Ω⊤
ij )−1 is attention-weighted precision aggregation that pulls agent
i’s precision toward its gauge-transported neighbors; and P
j (∂βij /∂Σi) DKL(qi ∥ Ωij qj ) is
gradient flow through the attention weights. In the general case (Eq. 35), αi = c0/(b0 +
DKL(qi∥pi)) depends on Σi through the trace and log-determinant terms of the full KL
divergence, so ∂αi/∂Σi̸ = 0. This gate gradient term vanishes only in the deterministic
limit (Eq. 74), where αi reduces to a function of μi alone. The coefficient −(1 + αi) on
Σ−1
i arises because the entropy of qi enters both the prior KL (weighted by αi) and each
neighbor KL (weighted by βij , summing to 1).
Since P
j ∂βij /∂Σi = 0 (the attention weights sum to unity), the attention weight cor-
rection vanishes whenever the KL divergences DKL(qi ∥ Ωij qj ) are uniform across neighbors.
It also vanishes in both the hard-attention (τ → 0, where βij concentrates on a single
neighbor) and uniform-attention (τ → ∞, where ∂βij /∂Σi ∝ 1/τ → 0) limits. Setting this
correction aside, the equilibrium condition ∂Fi/∂Σi = 0 yields the fixed-point equation:
Σ−1
i = 1
2

Σ−1
p,i + X
j
βij (Ωij Σj Ω⊤
ij )−1

 . (83)
Each agent’s precision is thus determined self-consistently as a β-weighted combination
of its prior precision and its neighbors’ transported precisions. When the state-dependent
prior coupling is weak (αi ≪ 1; cf. Section 3.7), the prior term becomes negligible. Simul-
taneously, the entropy contribution from the prior KL vanishes, changing the coefficient on
Σ−1
i from −2 to −1 in the gradient. The fixed point then reduces to Σ−1
i ≈ ⟨(Ωij Σj Ω⊤
ij )−1⟩β ,
the attention-weighted average of transported neighbor precisions. More generally, with
state-dependent precision αi (adaptive weight decay), the fixed-point interpolates as:
Σ−1
i = αi
1 + αi
Σ−1
p,i + 1
1 + αi
X
j
βij (Ωij Σj Ω⊤
ij )−1, (84)
a convex combination that continuously interpolates between prior-dominated (αi ≫ 1,
precision ≈ Σ−1
p,i ) and alignment-dominated (αi ≪ 1, precision ≈ ⟨(Ωij Σj Ω⊤
ij )−1⟩β ) regimes.
Since αi is itself state-dependent (Eq. 74), agents whose beliefs have drifted far from the prior
(large ∥μi − μprior∥) have small αi and therefore defer to their neighbors’ precisions, while
25
Dennis
agents near the prior retain strong prior anchoring. This then is the covariance analogue of
the adaptive weight decay described in Section 3.7. This fixed point is an attractor of the
gradient dynamics (Appendix B), so covariance alignment Σi ≈ Ωij Σj Ω⊤
ij emerges from the
variational dynamics itself rather than as an imposed constraint. This uncertainty channel
where precision propagates through the attention graph has no counterpart in standard
transformers and constitutes a concrete prediction of the general theory.
Gauge frame gradient. The gauge frame gradient ∂F/∂ϕi drives learning of the par-
allel transport operators Ωij = exp(ϕi) exp(−ϕj ). Its structure involves the differential of
the matrix exponential map and differs for compact (SO(N )) and general (GL(K)) gauge
groups; the full derivation is given in Appendix C. In practice, automatic differentiation
through the matrix exponential handles these gradients implicitly however we derive them
in detail.
4.3.3 Residual Connections as Variational Gradient Flow.
The residual (skip) connection is a defining architectural feature of modern transformers,
where each sublayer’s output is added to its input: xℓ+1 = xℓ + f (xℓ). This structure
emerges directly from the variational gradient flow. Identifying each transformer layer ℓ
with one step of natural-gradient descent on the free energy gives the Euler discretization:
μ(ℓ+1)
i = μ(ℓ)
i − ˜η Σi
∂F
∂μi μ(ℓ)
i
, (85)
where the Fisher preconditioning Σi and effective learning rate ˜η = ησ2 are as derived
above. This inherently has the form new = old + correction. The agent-based variational
model was derived from first principles, and this residual structure emerged as a consequence
of the gradient flow before the connection to transformer architecture was recognized. That
the same structure appears as a deliberate architectural choice in transformers (motivated by
gradient propagation considerations) suggests both frameworks capture the same underlying
computational principle. The backward gradient flow through the Jacobians ∂μ(ℓ+1)
j /∂μ(ℓ)
i
corresponds to variational message-passing, where each layer refines its beliefs based on
downstream prediction errors.
The strength of the identity path relative to the correction is governed by the state-
dependent precision αi (Section 3.7). When αi is large (agent near its prior), the identity
path dominates and the belief is strongly anchored; when αi is small (agent far from prior),
the correction terms dominate. In the constant-α limit, this reduces to a fixed-strength
residual connection with uniform regularization, as in standard transformers where the skip
connection weight is implicitly unity. The adaptive αi thus generalizes the fixed residual
connection to a state-dependent gating mechanism.
The full variational flow for all agent parameters takes the form:
θ(ℓ+1)
i = θ(ℓ)
i − ηθ ˜∇θi F θ(ℓ)
i
, θi ∈ {μi, Σi, ϕi}, (86)
where ˜∇θi denotes the natural gradient on each channel’s parameter manifold (Fisher-
Rao for μi and Σi; Killing form or pullback metric for ϕi, see Appendix C). Each channel
26
Attention and Transformers as Free Energy Minimization
possesses its own residual structure and effective learning rate. This multi-channel residual
flow corresponds to the residual stream interpretation of transformers (Elhage et al., 2021),
where each layer reads from and writes to a shared residual pathway. In the gauge-theoretic
framework, this residual stream is the belief state (μi, Σi, ϕi), accumulating incremental
updates from successive layers of variational refinement.
Layers as VFE iterations. The identification ℓ ↔ gradient step makes the notion of
“depth” a statement about how many steps of free energy minimization are performed.
A standard L-layer transformer corresponds to L natural-gradient steps, each with inde-
pendently parameterized attention weights β(ℓ)
ij (since W (ℓ)
Q , W (ℓ)
K , W (ℓ)
V differ across layers).
Within a single block one may also iterate the belief update T times with shared weights,
recomputing βij after each step; this gives T inner VFE iterations per layer and LT effec-
tive gradient steps in total. The standard transformer is the T = 1 case. The opposite
extreme, L = 1 with T → ∞, recovers a deep equilibrium model (Bai et al., 2019) in which
the forward pass iterates to a fixed point of the variational gradient. Intermediate regimes
interpolate between fast feedforward evaluation (T = 1, many layers) and slow iterative
inference (T ≫ 1, few layers), with the total depth of belief refinement governed by the
product LT .
4.4 Multi-Head Attention as Block-Diagonal GL(K) Structure
Standard transformer architectures employ multi-head attention, partitioning the dk-dimensional
embedding space into H independent heads (Vaswani et al., 2017):
μi = [h1
i , h2
i , . . . , hH
i ], ha
i ∈ Rdhead , dk = H × dhead. (87)
Each head computes attention independently using separate projection matrices (W a
Q, W a
K , W a
V ),
and the results are concatenated. The gauge-theoretic framework reveals this as implement-
ing a block-diagonal representation structure within the full GL(dk) gauge group.
4.4.1 Block-Diagonal Gauge Structure in GL(dk).
Recall from Section 2.1.3 that standard transformers implement GL(dk) gauge-theoretic
attention with gauge transport recovered via Ω = (σ2WK W ⊤
Q )−1 ∈ GL(dk). Multi-head
attention restricts this to a block-diagonal subgroup:
Ω =





Ω1 0 · · · 0
0 Ω2 · · · 0
... ... . . . ...
0 0 · · · ΩH




 , Ωa = (σ2W a
K (W a
Q)⊤)−1 ∈ GL(dhead), (88)
where each head a learns its own GL(dhead) gauge transformation. This block structure
means that components in different heads do not mix under gauge transport. Each head im-
plements an independent gauge-theoretic attention mechanism. Importantly, each per-head
transport Ωa ∈ GL(dhead) is full rank and invertible within its block, so the Mahalanobis
metric correction (Ωa)−⊤ from Eq. (42) remains well-defined. The rectangular shape of
the standard per-head projection matrices W a
Q ∈ Rdk ×dhead reflects the composition of (i)
27
Dennis
the orthogonal projection Pa : Rdk → Rdhead onto block a, with (ii) the GL(dhead) gauge
transport within that block. The rectangular outer shape is therefore the projection, not a
rank deficiency.
The full multi-head gauge group is thus the direct product:
Gmulti-head = GL(dhead)H ⊂ GL(dk), (89)
a (H · d2
head)-dimensional subgroup of the d2
k-dimensional GL(dk).
Per-head temperature. Since each head operates on an independent irrep block with
its own GL(dhead) gauge transport, the optimal temperature τa = κa
√dhead may differ
across heads. Heads encoding different types of structure (e.g., syntactic versus semantic,
local versus global) may benefit from sharper or broader attention distributions. In the
unreduced framework, each head has a learnable log-temperature log κa initialized from the
global κ and trained jointly with all other parameters. The per-head attention weights
become:
β(a)
ij = softmaxj − DKL(q(a)
i ∥ Ω(a)
ij q(a)
j )
κa
√dhead
!
, (90)
where q(a)
i denotes the belief projected onto head a’s subspace. This is a natural gener-
alization: just as the ALiBi slope ma is head-dependent (Section 3.4.5), the temperature κa
controls the sharpness of each head’s attention distribution independently. At initialization
(κa = κ for all a), this reduces to the standard global temperature. During training, heads
that develop sharp, peaked attention patterns (e.g., positional or syntactic heads) typically
learn smaller κa, while heads with diffuse contextual patterns learn larger values.
4.4.2 Geometric Interpretation via Lie Algebra Decomposition.
The Lie algebra gl(dk) has dimension d2
k, with generators spanning all dk ×dk matrices. The
block-diagonal restriction to gl(dhead)H projects out all off-diagonal blocks thereby retaining
only H · d2
head generators.
Therefore, each head’s Lie algebra gl(dhead) decomposes as:
gl(dhead) = sl(dhead) ⊕ R, (91)
where sl(dhead) consists of traceless matrices (volume-preserving transformations) and
R is the trace (overall scaling). This means each head independently controls the overall
magnitude of attention in that subspace (scaling) and how different dimensions within the
head interact (shearing and rotation).
4.4.3 Off-Diagonal Gauge Mixing: The Projected-Out Generators.
The block-diagonal restriction projects out a substantial portion of the full gauge algebra.
Decomposing gl(dk) into block-diagonal and off-diagonal components:
28
Attention and Transformers as Free Energy Minimization
gl(dk) =
HM
a=1
gl(dhead)
| {z }
block-diagonal (retained)
⊕ m|{z}
off-diagonal (projected out)
, (92)
the complement m has dimension d2
k −H ·d2
head. For typical architectures (e.g., dk = 512,
H = 8, dhead = 64), this is 5122 − 8 × 642 = 229,376 generators out of 262,144 total implying
that 87.5% of the gauge algebra is discarded by the multi-head factorization! These off-
diagonal generators are the inter-head mixing transformations. They are elements X ∈ m
that have non-zero entries only in off-diagonal blocks Xab (a̸ = b), each of which maps
Rdhead (head b) into Rdhead (head a).
Restoring these would yield a gauge transport with off-diagonal blocks:
Ωij =





Ω11
ij Ω12
ij · · · Ω1H
ij
Ω21
ij Ω22
ij · · · Ω2H
ij
... ... . . . ...
ΩH1
ij ΩH2
ij · · · ΩHH
ij




 ∈ GL(dk), (93)
where the off-diagonal blocks Ωab
ij ∈ Rdhead×dhead (a̸ = b) implement cross-head gauge
transport. When transporting the belief from token j to the frame of token i, head b’s
representation at j is rotated into head a’s subspace at i. The VFE alignment gradient
∂F/∂Ωab
ij for the off-diagonal blocks would then drive the system to learn these cross-head
correlations, encoding when one pattern appears in head b, it should be transported into
head a’s frame for comparison.
This is a qualitatively different mechanism from the standard output projection WO ∈
Rdk ×dk in transformers. The output projection mixes head outputs after attention. After
each head has independently computed its attention-weighted aggregation P
j βa
ij V a
j , the
results are concatenated and linearly mixed. Off-diagonal gauge transport, however, mixes
heads before attention: the cross-head blocks Ωab
ij enter the KL divergence that determines
the attention scores themselves. This makes the attention weights cross-head aware whereby
head a’s attention pattern can depend on the content of head b through the transport,
before any aggregation occurs at all. The full-GL(dk) transport is therefore more expressive
than the block-diagonal-plus-output-projection factorization, as it couples the attention
computation across heads rather than only the output computation.
From the perspective of the VFE gradient flow, the block-diagonal restriction constrains
the gradient ∇ϕF to lie within the block-diagonal sub-algebra L
a gl(dhead). Restoring the
off-diagonal components amounts to allowing the gradient to flow along the full gl(dk),
utilizing all d2
k generators. The projected-out gradient components ∇ϕF m represent cross-
head alignment forces that are present in the full VFE but zeroed out by the architectural
constraint. Whether these forces carry any useful signal is an empirical question with
significant implications.
4.5 Rotary Position Embeddings as Position-Dependent Gauge Frames
The prior-based positional mechanisms discussed in Section 3.4.5 (masking, ALiBi, relative
biases) all operate through the attention prior πj , entering as additive biases in the pre-
29
Dennis
softmax logits. A complementary class of positional encodings operates instead through
the gauge transport Ωij itself, making the gauge frame position-dependent. Rotary Position
Embeddings (RoPE) (Su et al., 2024) are the most widely adopted example. Here we derive
RoPE from our framework.
4.5.1 RoPE as a Gauge Subgroup Restriction.
RoPE applies position-dependent rotations to the query and key vectors before computing
the dot product:
Qi → R(θi) Qi, Kj → R(θj ) Kj , (94)
where R(θm) ∈ SO(2)dk /2 is a block-diagonal rotation matrix with dk/2 independent
2 × 2 rotation blocks:
R(θm) = diag
cos mθ1 − sin mθ1
sin mθ1 cos mθ1

, . . . ,
cos mθdk /2 − sin mθdk /2
sin mθdk /2 cos mθdk /2

, (95)
with frequencies θn = 10000−2n/dk . The attention logits become:
(R(θi)Qi)⊤(R(θj )Kj ) = Q⊤
i R(θi)⊤R(θj )Kj = Q⊤
i R(θj−i)Kj , (96)
where R(θj−i) = R(θi)⊤R(θj ) depends only on the relative position j − i.
In the gauge-theoretic framework, this has a straightforward interpretation. Recall from
Section 2.1.3 that WQW ⊤
K = σ−2Ω−⊤. With RoPE, the rotation is applied to query and
key vectors after projection, so the effective logit kernel becomes:
σ−2(Ω(RoPE)
ij )−⊤ = WQ R(θj−i) W ⊤
K , (97)
where R(θj−i) ∈ SO(2)dk /2 is sandwiched between the learned projections because it
acts in the projected dk-dimensional space. This partially relaxes Limit 2 (flat bundle):
the gauge connection acquires a position-dependent component, but it is restricted to the
compact (Abelian) subgroup SO(2)dk /2 ⊂ GL(dk).
4.5.2 Gauge-Theoretic Interpretation.
The factorization (97) identifies the position-dependent gauge frame as:
ϕ(pos)
i = diag(iθ1J, iθ2J, . . . , iθdk /2J) ∈ so(2)dk /2 ⊂ gl(dk), (98)
where J =  0 −1
1 0
 is the SO(2) generator and θn are the RoPE frequencies. The position-
dependent rotation is then R(θi) = exp(ϕ(pos)
i ), and the inter-agent rotation R(θj−i) =
exp(−ϕ(pos)
i ) exp(ϕ(pos)
j ) has exactly the form of a gauge transport operator (12) restricted
to the SO(2)dk /2 subgroup.
The multi-frequency structure of RoPE (θn = 10000−2n/dk ) corresponds to different
components of the gauge frame oscillating at different frequencies across positions: low-
frequency components (θn small, n large) encode coarse positional structure and remain
nearly aligned across distant positions, while high-frequency components (θn large, n small)
encode fine-grained local position and de-correlate rapidly.
30
Attention and Transformers as Free Energy Minimization
4.5.3 Interpolation Between Flat and Full Gauge.
RoPE occupies an intermediate position in the gauge-theoretic framework:
σ−2 exp(−ϕj )⊤ exp(ϕi)⊤
| {z }
Full GL(K) gauge (this work)
⊃ WQ R(θj−i) W ⊤
K
| {z }
RoPE
⊃ WQW ⊤
K
| {z }
Standard attention (flat bundle)
, (99)
where each expression is the logit kernel σ−2Ω−⊤
ij appearing in ˜sij = μ⊤
i (·)μj . Standard
attention uses a constant, learned gauge with WQW ⊤
K = σ−2Ω−⊤ (Limit 2 fully imposed).
RoPE augments this with a fixed, position-dependent SO(2)dk /2 rotation sandwiched be-
tween the projections (Limit 2 partially relaxed for positional information). The full gauge-
theoretic framework allows the complete transport Ωij to be agent-pair-dependent and
learned (Limit 2 fully relaxed). This hierarchy clarifies the design space such that posi-
tional encoding methods can be understood as choices about which subgroup of GL(K)
carries position-dependent structure, and which components remain position-independent
and learned from data.
An important distinction is that RoPE applies the position-dependent gauge only to the
attention computation (queries and keys) and not to the value aggregation. In our frame-
work, the full gauge transport Ωij mediates both the attention score DKL(qi∥Ωij qj ) and the
message aggregation mi = P
j βij Ωij μj . The asymmetry in RoPE with position-dependent
attention but position-independent values corresponds to factoring the gauge transport into
an attention gauge and a value gauge that need not coincide. This decomposition is sup-
ported (but not required) by the full framework.
4.6 Further Architectural Correspondences
Beyond the attention mechanism, several additional transformer components possess origins
within the gauge-theoretic framework. These are components that are derived from the
variational principle, components whose structural role is explained, and components that
are accommodated by the framework without being uniquely predicted.
4.6.1 Feed forward Networks and Non-linear Activations.
Standard transformers interleave attention sublayers with feedforward networks (FFNs),
typically two-layer MLPs with a nonlinear activation such as GeLU (Hendrycks and Gim-
pel, 2016). Within the gauge-theoretic framework, the FFN corresponds to the iterative
VFE dynamics that adjust beliefs between attention computations. The nonlinear acti-
vation function is not an independent architectural choice but a necessary consequence of
the information-geometric structure: variational inference on Gaussian beliefs produces a
Gated Linear Unit (GLU) (Dauphin et al., 2017) whose gate is a Boltzmann weight over
KL energies. The neural network is the computational substrate of a deeper information-
geometric dynamics; the activation functions used in practice (SiLU, GELU, ReLU) are
specific instantiations of this dynamics under successive limits.
31
Dennis
Message passing as a gated linear unit. Consider the message contribution of neigh-
bor j to the mean update (Eq. 80):
∆μi j = −˜η Σi βij (Ωij Σj Ω⊤
ij )−1μi − Ωij μj
. (100)
The message eij ≡ μi − Ωij μj is linear in the beliefs. The attention gate βij =
softmaxj
−DKL(qi∥Ωij qj )/τ  is manifestly a Boltzmann weight: each neighbor’s contri-
bution is weighted by exp(−Eij /τ )/Zi, where the energy Eij = DKL(qi∥Ωij qj ) depends on
the message magnitude through the Gaussian KL as:
DKL(qi∥Ωij qj ) = 1
2 e⊤
ij (Ωij Σj Ω⊤
ij )−1eij + (covariance terms). (101)
Each message’s contribution therefore has the general structure:
∆μi j ∝ eij
|{z}
linear
·
exp− 1
2τ ∥eij ∥2
Σ−1
j
+ · · · 
P
k exp− 1
2τ ∥eik∥2
Σ−1
k
+ · · · 
| {z }
Boltzmann gate
, (102)
which is a Gated Linear Unit: a linear function of the input gated by a Boltzmann distri-
bution over KL energies computed from the same input.
Binary limit: sigmoid as a special case. The general Boltzmann gate operates over all
ns sources simultaneously, where ns denotes the number of neighbors contributing to agent
i’s update (distinct from the sequence length N defined in Table 1; in the full attention
setting ns = N − 1, but the gate structure holds for any ns). In the special case where
only ns = 2 sources compete for influence on agent i, the Boltzmann distribution over two
states reduces to the logistic sigmoid σ, since exp(−E1/τ )/(exp(−E1/τ ) + exp(−E2/τ )) =
σ((E2 − E1)/τ ):
βi1 = σ
 DKL,i2 − DKL,i1
τ

, ∆μi 1 ∝ ei1 · σ
 ∥ei2∥2 − ∥ei1∥2
2σ2τ + · · ·

. (103)
This recovers the SiLU/Swish (Ramachandran et al., 2018) structure f (x) = x · σ(g(x)),
with the gate argument determined by relative KL energies. This binary case corresponds
to the FFN sublayer, where the hidden representation is split into two competing pathways
representing precisely two sources whose relative informational surprise determines the gate.
The sigmoid is thus not the general activation but instead is the ns = 2 limit of the
Boltzmann gate.
Boltzmann gate, Gaussian family, and connection to GELU. The information-
geometric origin of the gate determines its functional form. For example, in the isotropic
limit Σj = σ2I, the un-normalized gate for neighbor j is exp(−Eij /τ ) = exp(−∥eij ∥2/2σ2τ ),
a Gaussian function of the prediction error. The Boltzmann distribution over KL energies is
therefore a Gaussian mixture gate: each neighbor’s weight is a Gaussian in prediction-error
space, normalized by the partition function Zi = P
k exp(−∥eik∥2/2σ2τ ). The activation
32
Attention and Transformers as Free Energy Minimization
functions used in practice are specific instantiations of this information-geometric gate under
different limits:
SiLU(x) = x · σ(x) logistic sigmoid gate (ns = 2),
GELU(x) = x · Φ(x) Gaussian CDF gate, (104)
VFE(e) = e · exp(−∥e∥2/τ ′)/Z Boltzmann gate (general ns, Gaussian energy).
All three belong to the Gaussian family of gated linear units. The logistic sigmoid σ(x) =
1/(1+e−x) is the two-state Boltzmann distribution. The Gaussian CDF Φ(x) = 1
2 erfc(−x/√2)
is the integral of the Gaussian PDF ∝ exp(−x2/2) that appears as the VFE gate’s un-
normalized weight. These are members of the same exponential family, related by the
number of competing information sources (ns = 2 for sigmoid, ns > 2 for softmax) and the
functional form of the energy (E ∝ x for SiLU, E ∝ x2 for VFE).
ReLU is the zero-temperature (τ → 0) limit: the Boltzmann distribution concentrates on
the lowest-energy state and the gate becomes a step function Θ, giving x·Θ(x) = max(0, x).
Softmax-gradient correction. Beyond the GLU structure of the message-passing term
itself, the full VFE gradient contains an additional nonlinear correction from differentiating
βij with respect to μi:
∂βij
∂μi
= − βij
τ
"
∂DKL,ij
∂μi
− X
k
βik
∂DKL,ik
∂μi
#
, (105)
where DKL,ij ≡ DKL(qi∥Ωij qj ). This centered gradient (deviation of the per-source KL
gradient from its β-weighted mean) provides a higher-order non-linearity beyond the first-
order GLU gate. It creates positive feedback whereby similar beliefs yield higher β, which
in turn pulls beliefs closer together. For a single component, ∂βj /∂zj = βj (1 − βj ), which is
itself a sigmoid-shaped function of the logit. This correction is unique to the VFE setting,
where β depends on the same μ being updated. However, in standard transformers, the
analogous derivative appears only during backpropagation, not in the forward pass.
The same softmax-gradient non-linearity also appears in the covariance and gauge frame
channels. The Σi gradient (Eq. 82) contains the term P
j (∂βij /∂Σi) DKL(qi∥Ωij qj ), which
has the identical centered structure. Similarly, the gauge frame gradient (Eq. 148) yields
∂βij /∂ϕi = −(βij /τ )[∂Kij /∂ϕi − P
k βik ∂Kik/∂ϕi], the same self-normalizing form. Each
variational channel (μ, Σ, ϕ) thus contributes its own GLU-type activation and the compo-
sition of these channels within each VFE iteration produces a richer non-linear map than
any single channel alone.
FFN depth from VFE iterations. In the full VFE implementation, multiple iterations
of gradient descent with recomputed β at each step constitute the analogue of the FFN
sub-layer. Each iteration applies the GLU map (102) composed with the softmax-gradient
correction (105), producing increasingly complex nonlinear transformations. The number of
VFE iterations plays the role of FFN depth, as established in the “Layers as VFE iterations”
paragraph (Section 4.3.3).
33
Dennis
4.6.2 Token Embeddings as Prior Banks.
In standard transformers, token embeddings map discrete tokens to continuous vectors via
a learned embedding matrix Wembed ∈ RV ×d. In the gauge-theoretic framework, each token
type c ∈ {1, . . . , V } is associated with a learned prior distribution:
pc = N (μc
p, Σc
p), ϕc ∈ gl(K), (106)
comprising a prior mean μc
p, prior covariance Σc
p, and gauge frame coefficients ϕc. At the
start of each sequence, agent i’s belief is initialized to the prior of its token type: qi(0) =
ptoken(i). The collection {(μc
p, Σc
p, ϕc)}V
c=1 constitutes a prior bank, which is the generalization
of the embedding matrix. In the deterministic, isotropic limit (Σc
p → σ2I, ϕc absorbed), only
the prior means survive and the prior bank reduces to Wembed = [μ1
p, . . . , μV
p ]⊤, recovering
the standard embedding matrix.
4.7 Summary
Under the successive limits: (i) isotropic beliefs with σ−2 absorbed into learned projec-
tions, (ii) constant gauge (Ωij = Ω for all pairs), and (iii) learned gauge via projec-
tions (WQW ⊤
K = σ−2Ω−⊤ ∈ GL(dk)), variational natural-gradient descent on the gauge-
equivariant free energy is:
θk+1 = θk − η ˜∇θF(θk) (107)
and recovers the standard training update:
θk+1 = θk − η ˜∇θL(θk), (108)
under the identifications: F (free energy) ↔ L (loss + regularization); attention weights
βij ↔ transformer attention αij ; observations − log p(o | μ) ↔ loss function; and symmetry
breaking ↔ training.
Beyond the core attention formula, the gauge-theoretic framework provides a unified
account of the full transformer architecture (Table 3). Causal masking, ALiBi, relative
position biases, and sliding window attention all emerge as special cases of the non-uniform
attention prior πk in the mixture-of-sources model (Section 3.4.5). Rotary position em-
beddings correspond to position-dependent gauge frames restricted to SO(2)dk /2, partially
relaxing the flat-bundle limit while retaining the learned constant gauge (Section 4.5). Layer
normalization is a projection onto the sl(K) subalgebra that eliminates key-norm bias at
its algebraic origin. Residual connections are the inherent structure of iterative variational
gradient descent (Section 4.3.3), with the state-dependent precision αi generalizing the fixed
skip connection to a context-dependent gate. The FFN sublayer corresponds to multiple
VFE iterations with recomputed attention weights, and the message-passing term derives a
Boltzmann-gated linear unit (Eq. 102) in the same exponential family as GELU and SiLU
(Section 4.6.1). Weight decay is the prior coupling αDKL(qi∥pi) in the deterministic limit.
The GL(K) gauge invariance theorem (Theorem 1) is essential as it guarantees that the
learned projections WQ, WK are valid gauge transformations, explaining why transformers
can learn arbitrary (invertible) attention patterns. Each row of Table 3 represents a degree
of freedom that the full gauge-theoretic framework retains and that standard transformers
34
Attention and Transformers as Free Energy Minimization
collapse into a simpler form. The framework thus provides not only an explanatory account
of existing architectures but a taxonomy of generalizations, indexed by which simplifying
limits are relaxed.
5 Simulations and Empirical Validation
5.1 Experimental Design
Hardware details are provided in Appendix E.1.
5.1.1 Transformer Validation Protocol
To empirically validate the equivalence between our gauge attention rule and standard
transformer self-attention, we conducted a quantitative comparison using a pretrained
BERT-base-uncased model from HuggingFace Transformers (Devlin et al., 2018). We con-
structed a corpus of 105 diverse English passages spanning multiple genres, registers, and
subject domains, tokenized each using BERT’s WordPiece tokenizer, and performed full
forward passes while extracting hidden states from all 12 layers.
The validation proceeds in seven phases. All primary analyses use the theory-predicted
temperature τ = 19.0 ≈ 2√d for d = 64 (the value of τ that our gauge framework pre-
dicts from dimensional scaling; see Section 5.2.2): (1) single-passage analysis establishing
per-head agreement; (2) multi-passage corpus validation across all 105 passages, quantify-
ing cross-passage stability via standard errors and confidence intervals for each of the 144
attention heads; (3) a temperature sweep across the full corpus at 10 values of τ from 0.1 to
50, confirming the empirical optimum at τ = 19.0 and its consistency with the theoretical
prediction τopt = 2√d; (4) multi-model validation across five architectures (BERT-base,
BERT-large, DistilBERT, RoBERTa, ALBERT); (5) mathematical identity decomposition
verifying Q · K/√d = (−∥Q − K∥2 + ∥Q∥2 + ∥K∥2)/2√d to machine precision; (6) at-
tention entropy analysis confirming that H(β) ≈ H(α) at the optimal temperature; and
(7) sequence-length sensitivity analysis from N = 16 to N = 512 tokens.
For each layer L and head H, we extracted query, key, and value matrices
Q(L,H), K(L,H), V (L,H) ∈ RT ×d,
where T is the token sequence length and d = 64 is the head dimension. We compared two
attention mechanisms:
Standard transformer attention:
αij = softmaxj
 Qi · Kj
√d

(109)
KL gauge attention (flat bundle):
β(flat)
ij = softmaxj

− ∥Qi − Kj ∥2
τ

(110)
For each (L, H) pair across all 144 attention heads and each passage, we computed
the Pearson correlation between corresponding attention rows α and β, then aggregated
per-head correlations across the 105-passage corpus to obtain means, standard errors, and
35
Dennis
95% confidence intervals. We additionally computed the fraction of tokens with identical
argmax attention, arg maxj αij = arg maxj βij (peak match), and the key-norm correlation
ρ(∥Kj ∥2, attention weight) for both α and β on each passage.
5.1.2 GL(K) Language Modeling
To validate the complete gauge-theoretic framework beyond attention mechanism corre-
spondence, we trained a GL(K) gauge transformer on WikiText-103 language modeling
(102M tokens, vocabulary size 50,257). These experiments test whether variational free
energy minimization can drive actual language learning, with the gauge structure (μ, Σ, ϕ)
evolving jointly during training. The gauge VFE architecture contains no MLPs, pointwise
activation functions (ReLU, GELU, etc.), or learned attention projections (WQ, WK , WV );
only a linear output projection to vocabulary logits is retained. The model does employ
nonlinear operations inherent to the framework: softmax for attention normalization and
matrix exponentials for gauge transport. All other computation derives from KL diver-
gences, parallel transport via Ωij = exp(ϕi) exp(−ϕj ), and natural gradient descent on the
Fisher-Rao manifold. No positional encoding was used. The causal attention mask provides
implicit sequential structure by restricting each token’s context to preceding positions, and
the gauge transport operators Ωij can in principle learn position-dependent transforma-
tions through the token-specific gauge frames ϕi. We note that this design choice means
the model cannot distinguish absolute positions within the causal window (only relative
ordering is implicit in the mask boundary), which may limit performance compared to ar-
chitectures with explicit positional encoding. The standard transformer baseline likewise
uses no positional encoding, ensuring a controlled comparison.
We report results for a GL(10) gauge VFE model with embedding dimension K = 90,
decomposed into 9 fundamental irreps of dimension 10 (yielding 9 attention heads), context
length N = 128, batch size 16, and 50,000 training steps (58.8M parameters). The model
uses VFE-dynamic mode with evolving covariance (Σ) and gauge frame (ϕ), full (non-
diagonal) covariance matrices, and learnable per-component learning rates. A standard
transformer with dot-product attention and MLP sublayers (17.6M parameters) trained
on the same data serves as a baseline. Full configurations are documented in the code
repository.
Variational Inference Structure (E-step/M-step). We utilize an expectation-maximization
structure as follows:
E-step (Belief Inference): Given current parameters, update agent beliefs {qi} via
natural gradient descent on the free energy:
q(t+1)
i ← q(t)
i − ηE ˜∇qi F qi=q(t)
i
, (111)
where ˜∇ denotes natural gradients projected via the Fisher-Rao metric. This corre-
sponds to variational inference within the forward pass.
M-step (Learning): After belief inference, gradients flow via backpropagation through
the entire computation graph:
θ ← θ − ηM ∇θL[{q∗
i }, θ], (112)
36
Attention and Transformers as Free Energy Minimization
Algorithm 1 GL(K) Gauge Transformer Training Loop
Require: Token sequence (x1, . . . , xN ); learnable parameters θ =
{μembed, Σembed, ϕembed, Wout}
1: Initialize beliefs from priors: μi ← μxi
embed, Σi ← Σxi
embed, ϕi ← ϕxi
embed
2:
3: — E-step: Variational belief inference (within forward pass) —
4: for t = 1 to Tinner do
5: Compute transport: Ωij ← exp(ϕi) exp(−ϕj ) ▷ Gauge parallel transport
6: for each head h (irrep block) do
7: D(h)
ij ← DKL(q(h)
i ∥ Ω(h)
ij q(h)
j ) ▷ Pairwise KL divergences
8: β(h)
ij ← softmaxj

− D(h)
ij / κh

▷ KL attention weights
9: end for
10: — Mean gradient (direct + softmax nonlinearity, Eq. 105):
11: ∇μF ← α Σ−1
p (μi − μp)
| {z }
prior
+ P
h
h X
j
β(h)
ij ∇μD(h)
ij
| {z }
direct alignment
+ X
j
D(h)
ij
∂β(h)
ij
∂μi
| {z }
softmax nonlinearity
i
+ ∇μLobs
| {z }
observation
12: ˜∇μ ← Σi · ∇μF ▷ Natural gradient (Fisher preconditioning)
13: μi ← μi − ηE ˜∇μ ▷ Trust-region-constrained update
14: — Covariance gradient (Eq. 82):
15: ∇ΣF ← α
2 (Σ−1
p − Σ−1
i ) + 1
2
P
h
P
j β(h)
ij
(Ω(h)
ij Σ(h)
j Ω(h)⊤
ij )−1 − Σ(h)−1
i

16: Σi ← RetractSPD(Σi, −ηΣ Σi ∇ΣF Σi) ▷ SPD retraction with Fisher metric
17: — Gauge frame gradient (differentiating through β(ϕ), Eq. 146):
18: ∇ϕF ← ∇ϕi
h P
h
P
j β(h)
ij (ϕ) D(h)
ij (ϕ) + αϕ
2 ∥ϕi∥2
i
▷ Autograd through β
19: ϕi ← Retractg(ϕi, −ηϕ ∇ϕF) ▷ Lie algebra retraction
20: end for
21:
22: — M-step: Parameter learning (backpropagation) —
23: Compute loss: L = LCE(μ∗
i , Wout) + ˆα P
i DKL(q∗
i ∥ pi) + αϕ
2 ∥ϕ∥2
24: θ ← θ − ηM ∇θL ▷ Adam on all embeddings and Wout
where θ = {μembed, Σembed, ϕembed, Wout}. Within each E-step iteration, belief gradients
are computed analytically from detached snapshots of μ, Σ, and ϕ, so the VFE dynamics are
not differentiated through via autograd. However, the final beliefs {q∗
i } retain a gradient
path to the initial embeddings through the additive update chain μ(t+1) = μ(t) + ∆μ(t),
acting as a straight-through estimator: the M-step learns embeddings whose initializations
lead to good converged beliefs, without requiring second-order differentiation through the
E-step trajectory.
5.1.3 Belief and Prior Initialization
Agent beliefs qi = N (μi, Σi) represent context-dependent semantic interpretations that
evolve through communication. Token priors pc = N (μc
p, Σc
p) and gauge frame coefficients
37
Dennis
ϕa
c ∈ R are initialized from N (0, σ2
init) with σinit = 0.3, serving as learnable token represen-
tations analogous to embedding matrices.
For GL(K), the gauge frame ϕc directly parameterizes elements of gl(K), with trans-
port operators Ωi = exp(ϕi) ∈ GL+(K) computed via matrix exponential. Since the matrix
exponential exp : gl(K) → GL+(K) is surjective onto the identity component, this pa-
rameterization covers all orientation-preserving transport operators. At the start of each
training sequence, agent beliefs are initialized to match their token priors:
qi(t = 0) = ptoken(i), (113)
where token(i) denotes the token type at position i. Thus μi(0) = μtoken(i)
p and Σi(0) =
Σtoken(i)
p . All instances of the same token begin with identical beliefs regardless of position;
context-dependent representations emerge through belief evolution during the E-step.
5.2 Empirical Results
5.2.1 Overall Agreement and Temperature Optimization
We validated the gauge-aligned KL attention against standard dot-product attention across
all 144 heads of BERT-base-uncased on a corpus of 105 diverse passages, using τ = 19.0 as
the default temperature throughout. This value is not a free parameter: it is the theory-
predicted optimum τ ≈ 2√d for head dimension d = 64 (confirmed by the temperature
sweep in Section 5.2.2).
At τ = 19.0, the single-passage analysis yields mean correlation ¯r = 0.795 ± 0.229 and
median ˜r = 0.880 across all 144 heads, with 94 heads (65.3%) exceeding r > 0.8 and 67
heads (46.5%) exceeding r > 0.9. All 144 heads achieve p < 0.001 significance.
The multi-passage corpus validation across all 105 passages yields grand mean ¯r =
0.804 ± 0.182 and grand median ˜r = 0.876, with 93 of 144 heads (64.6%) exceeding mean
r > 0.8 and 59 heads (41.0%) exceeding mean r > 0.9. The range of per-head mean
correlations spans [0.049, 0.999].
Cross-passage stability. The robustness of these results across the 105-passage corpus
is a central finding. Standard errors for each head’s correlation are uniformly small: mean
SE = 0.004, median SE = 0.004, and maximum SE = 0.016, with all 144 heads achieving
SE < 0.05 (Figure 4). This demonstrates that the α–β agreement is a robust property of
BERT’s learned representations, not an artifact of any particular input passage.
The distribution of per-head correlations (Figure 3) reveals substantial heterogeneity
consistent with functional specialization. Some heads achieve near-perfect agreement (e.g.,
Layer 2, Head 0: mean r = 0.999, SE = 0.0004), while a minority of heads show weaker
correspondence (r < 0.3), typically those implementing positional or syntactic attention
patterns orthogonal to the KL-based distance metric. All 144 heads achieve p < 0.001
significance on 100% of passages, confirming that the agreement is statistically robust across
the full head population.
Figure ?? shows the per-head correlation structure across layers. Deeper layers (8–11)
tend to exhibit higher correlations with KL attention, consistent with the expectation that
semantic representations become more structured and therefore more amenable to distance-
based comparison in deeper layers.
38
Attention and Transformers as Free Energy Minimization
Figure 2: Temperature sweep across 105 passages. Grand mean Pearson correlation
between transformer and KL attention weights as a function of temperature τ , averaged
over all 144 heads and 105 passages. Shaded regions indicate 95% confidence intervals. The
empirical optimum at τ = 19.0 (¯r = 0.804, 95% CI [0.771, 0.838]) is consistent with the
theoretical prediction τopt = 2√d = 16 for d = 64. Correlation increases monotonically
from τ = 0.1 (¯r = 0.654) to τ = 19, then decreases at higher temperatures as the softmax
becomes overly diffuse.
Scope of the Validation. An important caveat is in order. The high correlation be-
tween KL-based attention βij ∝ exp(−∥Qi − Kj ∥2/τ ) and dot-product attention αij ∝
exp(QiK⊤
j /√dk) follows in part from an algebraic identity: when key norms ∥Kj ∥ are ap-
proximately constant (as enforced by layer normalization in BERT), the squared Euclidean
distance ∥Qi −Kj ∥2 = ∥Qi∥2 +∥Kj ∥2 −2QiK⊤
j differs from −2QiK⊤
j only by terms that are
approximately constant under softmax. The non-trivial content of this validation is there-
fore not the correlation per se (that much is guaranteed by the algebra) but rather four
quantitative predictions that the identity alone does not supply: (a) the optimal tempera-
ture τ ≈ 2√dk, confirmed empirically at τ = 19.0 (within 19% of the prediction 2√64 = 16),
whose value is fixed by the dimensional scaling of dot-product variance in the gauge frame-
work; (b) the existence, sign, and magnitude of a residual key-norm bias, a prediction
specific to the gauge-theoretic decomposition that has no counterpart in the algebraic iden-
tity; (c) the entropy matching H(β) ≈ H(α) at the optimal temperature (Section 5.2.4),
confirming that the temperature does not merely maximize correlation but recovers the
correct distributional shape; and (d) the per-layer variation in correlation strength, which
reflects representational structure learned during training and is not predicted by the math-
ematical identity alone. The multi-passage design (105 passages, all heads, full temperature
sweep) ensures that these findings are not artifacts of any particular input; the cross-passage
standard errors (SE < 0.02 for all 144 heads) confirm reproducibility.
5.2.2 Temperature Scaling and Finite-Dimensional Corrections
The multi-passage temperature sweep independently confirms that τ = 19.0 maximizes
the grand mean correlation across the full 105-passage corpus, yielding ¯r = 0.804 (95%
CI [0.771, 0.838]). This represents a 19% deviation from the theoretical prediction τopt =
2√d = 16 for d = 64. This does not represent a failure of theory but is a manifestation of
finite-dimensional corrections that our framework explicitly predicts.
39
Dennis
Figure 3: Distribution of per-head correlations across 105 passages. Each head’s
correlation is its mean r(α, β) across the full 105-passage corpus at the theory-predicted
optimal temperature τ = 19.0. The grand mean is ¯r = 0.804 and grand median is ˜r = 0.876;
93 of 144 heads (64.6%) exceed r > 0.8 and 59 heads (41.0%) exceed r > 0.9. The rightward
skew reflects that the majority of heads show strong agreement, with a tail of weakly
correlated heads exhibiting specialized (e.g., positional) attention patterns.
Temperature scaling emerges from competing effects: (1) dot products QiK⊤
j scale
as O(dk) in magnitude, while (2) key-norm fluctuations scale as O(√dk) under high-
dimensional measure. The factor √dk normalizes pre-softmax logits to O(1), the appropri-
ate scale for softmax attention. However, at finite dimensions, subdominant fluctuations
O(√d) induce corrections of order unity to the optimal temperature.
We validated these predictions by computing the coefficient of variation (CV) of key
norms, which directly measures incomplete bias cancellation. For d = 64, theory predicts
CV = p2/d = 0.177 (17.7%), representing fundamental O(1/√d) fluctuations. Monte
Carlo simulations with learned projection matrices yield CV = 0.240 ± 0.001 (24.0%),
where the amplification reflects typical BERT-like architectures.
40
Attention and Transformers as Free Energy Minimization
Figure 4: Cross-passage stability of per-head correlations. Standard errors across
the 105-passage corpus at τ = 19.0 are uniformly small (mean SE = 0.004, median = 0.004,
max = 0.016), with all 144 heads achieving SE < 0.05. This demonstrates that the α–β
agreement is a reproducible property of BERT’s learned representations rather than an
artifact of specific input text.
This is consistent with both observed phenomena: the temperature shift from τ = 16 to
τ = 19 follows from the embedding scale, with ratio τemp/τtheory = 1.188 closely tracking the
CV = 17.7% predicted by dimensional scaling (the 18.75% empirical deviation lies within
the range expected from learned-weight amplification); and the magnitude of key-norm
bias (discussed below) falls within the predicted range from 24% norm heterogeneity. The
multi-passage confidence interval [0.771, 0.838] at the optimum confirms that these results
are stable across diverse inputs.
Both effects are quantitative validations of the finite-dimensional analysis, with the
gauge framework correctly predicting their existence and magnitude from dimensional scal-
ing alone.
The temperature parameter τ plays a dual role. Geometrically, it represents the inverse
stiffness of gauge alignment: higher τ allows greater tolerance for misalignment, softening
attention peaks and distributing weight more uniformly among agents. Statistically, τ =
41
Dennis
2√dk reflects both the dimensional normalization √dk (encoding the belief precision σ2 =√dk from concentration of measure) and the quadratic coefficient 1
2 in the KL divergence
between isotropic Gaussians. The broad plateau from τ = 15 to τ = 25 (Figure 2) indicates
that the optimum is not sharply tuned, consistent with the theoretical picture of a smooth
dimensional correction rather than a resonance.
5.2.3 Key-Norm Bias and Layer Normalization
A central prediction of our gauge-theoretic framework is the emergence of a key-dependent
bias term that modulates attention beyond simple query-key compatibility. Recall from
Section 4.2 that prior to absorbing σ−2 into the learned projections, the full KL-derived
attention logit contains a term − 1
2σ2 ∥μj ∥2 depending only on the key embedding norm.
After the absorption (Section 4.2.2), this bias term becomes implicit in the learned weight
matrices. However, the bias is observable through its effect on the relationship between key
norms and attention weights. In the pre-absorption form, the attention score is:
β(flat)
ij = softmaxj

− ∥Qi − Kj ∥2
τ − λ∥Kj ∥2

, (114)
where λ > 0 is the residual bias coefficient (equal to 1/(2σ2) before absorption, and
implicitly encoded in the learned projections afterward). The term −λ∥Kj ∥2 represents
an intrinsic salience depending only on key vector norm. This bias is gauge-geometric in
origin: each key carries information not only about semantic content but also about its
representation magnitude in the embedding space.
Our theory predicts this bias should approximately cancel. As shown in Section 4.2.1,
the Mahalanobis metric (42) already accounts for the gauge transport’s effect on norms, so
the residual bias depends on the raw embedding norms ∥μj ∥2 regardless of the gauge group.
High-dimensional concentration then implies: for embeddings μ(i)
j ∼ N (0, σ2
0 ) in Rdk (where
σ2
0 denotes the per-component embedding variance),
∥μj ∥2 =
dkX
i=1
(μ(i)
j )2 = dkσ2
0 ± O(σ2
0
pdk), (115)
with relative fluctuations O(1/√dk) → 0 as dk → ∞. Thus ∥μj ∥2 ≈ dkσ2
0 becomes
approximately constant, and the bias term becomes approximately constant in j, cancelling
under softmax normalization. At finite dimensions, this cancellation is incomplete, leading
to a residual per-key bias.
The multi-passage analysis confirms this effect (Figure 5). For dot-product attention
(α), the mean absolute correlation between key norms ∥Kj ∥2 and per-token attention weight
is |ρ| = 0.139 across all 144 heads, while for KL-based attention (β) at τ = 19.0, the
corresponding value is |ρ| = 0.256. Both metrics aggregate across the full 105-passage
corpus, providing robust cross-passage estimates.
The stronger key-norm effect for β at τ = 19 compared to the |ρ| = 0.164 observed at
τ = 1 is itself a theoretical prediction: at higher temperature, the softmax distribution is
softer, so the norm-dependent bias term −λ∥Kj ∥2 contributes a larger relative perturbation
42
Attention and Transformers as Free Energy Minimization
Figure 5: Key-norm bias analysis across 105 passages. Per-head correlations
ρ(∥Kj ∥2, attention weight) for both dot-product attention (α) and KL-based attention (β),
aggregated across the full corpus. The predominantly negative correlations for β confirm
the theoretical prediction that the KL-derived bias term −λ∥Kj ∥2 suppresses attention to
high-norm keys.
to the flattened attention distribution. At τ = 1, β is extremely peaked (the nearest-
neighbor dominates regardless of norms), masking the bias. At τ = 19, the bias has room
to express.
Although the corpus-averaged magnitude is modest, individual heads show strong key-
norm effects (e.g., Layer 1, Head 6 has ρα = 0.889, reflecting a head where key-norm
variation dominates the attention pattern), and the cross-passage per-head key-norm corre-
lations (Figure 6) confirm that the sign and magnitude of the bias are reproducible across
inputs.
The predominantly negative sign for KL-based attention (β) is precisely what the the-
ory predicts: the key-norm bias term −λ∥Kj ∥2 contributes negatively to attention logits,
suppressing attention to high-norm keys. The weaker and sign-mixed correlations for dot-
product attention (α) reflect the query-dependent cancellation of key-norm terms under
softmax, which is exact only when ∥Kj ∥ is constant across tokens.
These findings suggest why norm-regulating mechanisms are prevalent in transformer
architectures. Layer normalization, RMSNorm, and related techniques enforce approxi-
mately constant norms across tokens, satisfying the gauge-theoretic cancellation condition
that should hold asymptotically in high dimensions. Without such regulation, key-norm
heterogeneity introduces systematic bias that modulates effective attention allocation. Our
43
Dennis
Figure 6: Cross-passage stability of key-norm bias. Per-head key-norm correlation
distributions across 105 passages, demonstrating that the sign and magnitude of the key-
norm bias are reproducible across diverse inputs for both α and β attention mechanisms.
gauge theory derives the condition (∥μj ∥2 ≈ C) but does not uniquely predict which normal-
ization scheme satisfies it; rather, it explains why some form of norm regulation is necessary:
transformers approximate variational inference in a gauge geometry, and proper inference
requires frame-independent comparisons achieved only when key norms are regulated.
5.2.4 Attention Entropy Matching
A necessary condition for the gauge-theoretic interpretation is that the KL-distance atten-
tion β not only correlates with dot-product attention α but reproduces its distributional
shape: that is, the entropy of the two attention distributions should match. If τ merely
maximized correlation by collapsing β toward a delta function (low entropy) or inflating it
toward uniformity (high entropy), the agreement would be artifactual.
We computed the Shannon entropy H(·) = − P
j pj log pj of both α and β for each
(layer, head) pair, averaged across 20 passages from the corpus. At the optimal temperature
τ = 19.0, the entropies are closely matched:
H(α) = 1.774 nats, H(β) = 1.784 nats, H(β)/H(α) = 1.076. (116)
The entropy ratio H(β)/H(α) ≈ 1.08 confirms that β at τ = 19 is neither collapsed nor
diffuse relative to α: the two distributions occupy the same region of the probability simplex.
This is a non-trivial consequence of the dimensional scaling. At τ = 1 (no rescaling), the
squared distances ∥Qi −Kj ∥2 are unscaled, producing extremely peaked β distributions with
H(β) ≪ H(α). The identity QiK⊤
j /√d = (−∥Qi − Kj ∥2 + ∥Qi∥2 + ∥Kj ∥2)/2√d predicts
that the natural scale for the distance-based logits involves a factor of 2√d; the empirical
optimum τ = 19 ≈ 2√64 = 16 recovers this scale, and the entropy matching confirms it
(Figure 7).
The layer-depth structure of entropy is also informative. Early layers exhibit higher
attention entropy (H(α)early = 2.225 nats for layers 0–2) than late layers (H(α)late = 1.801
nats for layers 9–11), reflecting a transition from diffuse, exploratory attention in early
layers to focused, task-specific attention in deeper layers. Both α and β track this pattern,
with the entropy ratio remaining near unity across all depths.
44
Attention and Transformers as Free Energy Minimization
Figure 7: Attention entropy by layer depth. Shannon entropy H(α) and H(β) for dot-
product and KL-distance attention, averaged across heads and 20 passages at τ = 19.0. The
close match H(β)/H(α) = 1.076 confirms that the optimal temperature recovers not just
correlation but distributional shape. Early layers exhibit higher entropy (diffuse attention)
while later layers are more peaked (selective attention), a pattern tracked by both α and β.
5.2.5 Multi-Model Generalization
To test whether the α–β correspondence generalizes beyond BERT-base, we replicated the
multi-passage analysis (20 passages, τ = 19.0) across five pretrained transformer archi-
tectures spanning different model families, depths, and parameter counts. All models use
dhead = 64, so the same temperature τ = 19.0 applies.
The correspondence holds across all five architectures (¯r > 0.6 in every case), with
ALBERT achieving the highest agreement (¯r = 0.851). ALBERT’s parameter-sharing
across layers forces all layers to use the same projection matrices, which may regularize
the attention patterns toward the distance-based structure predicted by the gauge theory.
RoBERTa’s lower correlation (¯r = 0.612) is consistent with its different pretraining objec-
tive (no next-sentence prediction, dynamic masking, larger batches), which may encourage
attention patterns that rely more heavily on the norm-bias terms that the flat-bundle ap-
proximation neglects.
45
Dennis
The theoretical prediction τ ≈ 2√dhead holds uniformly: all five models share dhead = 64,
and all achieve their best correspondence at τ = 19.0. This confirms that the optimal tem-
perature is a function of head dimension, not of model family, depth, or training procedure.
5.2.6 Sequence-Length Sensitivity
To verify that the α-β correspondence is not an artifact of a particular sequence length, we
evaluated the correlation across six sequence lengths from N = 16 to N = 512 tokens at
τ = 19.0 (Table 5).
The mean correlation remains robust across all sequence lengths (¯r > 0.71), with a mild
monotonic decrease from ¯r = 0.798 at N = 32 to ¯r = 0.714 at N = 512. This degradation
is consistent with the key-norm bias analysis (Section 5.2.3): longer sequences contain more
diverse token types, increasing key-norm variance and amplifying the residual bias term that
the flat-bundle isotropic limit neglects. The stability across a 32× range of sequence lengths
confirms that the correspondence is a structural property of the attention mechanism, not
an artifact of sequence statistics.
5.3 GL(K) Language Modeling: The Full General Model
Unlike the BERT validation above, which tested the degenerate limits against an existing
standard transformer across 105 diverse passages, our language modeling experiments imple-
ment the most general forms of the gauge-theoretic framework. These GL(K) models retain
full non-isotropic covariances Σi, non-trivial gauge transport Ωij = exp(ϕi) exp(−ϕj ), and
KL-divergence attention. This is a genuine test of whether the gauge-theoretic structure,
in its unreduced form, can implement meaningful language modeling.
5.3.1 Training Dynamics
Figure 8 shows the training dynamics for our best GL(K) gauge VFE model: a GL(10) con-
figuration with K = 90 (9 heads of dimension 10). The model achieves stable convergence
with perplexity dropping from 50,257 at initialization to ∼75 at convergence; a ∼671× im-
provement over random chance (vocabulary size 50,257). The gradient norms for μ (belief
means), Σ (covariances), and ϕ (gauge frames) evolve with distinct dynamics; the gauge
frame gradients (ϕ, red) exhibit characteristic oscillations as the model learns optimal frame
alignments. Training and validation curves track closely throughout training, indicating the
gauge-theoretic inductive bias provides effective regularization.
Figure 9 shows the learned KL-divergence attention patterns at convergence for four of
the nine heads (H0, H2, H4, H8). Despite having no learned attention projections (WQ, WK ,
WV ), the gauge VFE model learns structured, non-uniform, and head-specialized attention
patterns purely through gauge transport and KL-divergence dynamics. All four heads ex-
hibit the causal mask and strong early-position attention (vertical stripes at key positions 0–
10), reminiscent of the “attention sink” phenomenon observed in standard transformers (?).
Beyond these shared features, the heads develop qualitatively distinct attention strategies.
Head 0 (Figure 9a) distributes attention broadly and diffusely across the available context,
consistent with a global information-gathering role. Head 2 (Figure 9b) develops sharper
vertical stripes at specific key positions corresponding to high-frequency or structurally im-
portant tokens, suggesting a token-selective pattern. Head 4 (Figure 9c) exhibits prominent
46
Attention and Transformers as Free Energy Minimization
(a) Generalization gap (val BPC − train BPC)
for the GL(10), K = 90 gauge VFE model.
(b) Training curves for the GL(10), K = 90
gauge VFE model: loss, perplexity, bits per to-
ken, and gradient norms.
Figure 8: Training dynamics for the GL(10) gauge VFE model with K = 90 embedding
dimensions (9 heads of dimension 10) on WikiText-103. (a) The generalization gap (val-
idation BPC minus training BPC) remains small throughout training, indicating effective
regularization from the gauge-theoretic inductive bias. (b) Training curves showing loss,
perplexity, bits per token, and gradient norms; perplexity drops from >900 to ∼75 over
60,000 steps.
horizontal bands at specific query positions—likely punctuation or sentence boundaries—
where the query attends broadly to its entire available context, consistent with a positional
or syntactic aggregation role. Head 8 (Figure 9d) concentrates attention in a narrow band
near the diagonal, exhibiting a locality bias reminiscent of the sliding-window attention
prior (Eq. 26), but learned rather than imposed. This functional diversification—broad
context, token-selective, query-driven aggregation, and local recency—emerges solely from
the KL-divergence geometry on the statistical manifold, without any architectural inductive
bias toward head specialization. The log-scale visualization reveals that attention weight
variation spans approximately five orders of magnitude (10−5 to 100), confirming that the
model learns highly selective attention from the gauge-theoretic dynamics alone.
5.3.2 Results
Table 6 presents our best GL(K) gauge VFE model alongside a standard transformer base-
line trained on the same data.
The gauge VFE model achieves test perplexity 74.9 (∼671× improvement over random
chance, i.e., vocabulary size 50,257), demonstrating that the unreduced gauge-theoretic
framework, with KL-divergence attention, gauge transport, and natural gradient dynamics
on the statistical manifold, is sufficient for non-trivial language modeling. The standard
transformer baseline achieves test PPL 65.0, approximately 1.15× lower, reflecting the com-
putational advantages of dot-product attention with learned projections and MLP sublayers.
47
Dennis
(a) Head 0: broad, diffuse contextual attention
distributed across the full available context.
(b) Head 2: token-selective attention with ver-
tical stripes at specific high-frequency key posi-
tions.
(c) Head 4: query-driven aggregation with
prominent horizontal bands at structurally
salient positions.
(d) Head 8: local attention concentrated near
the diagonal, exhibiting a learned recency bias.
Figure 9: Learned per-head KL-divergence attention patterns (log10 βij ) for four of nine
heads in the GL(10) gauge VFE model (K = 90) at convergence (step 60,000) on a repre-
sentative WikiText-103 passage. All heads share the causal mask and early-position atten-
tion sinks, but develop qualitatively distinct strategies—broad context (a), token-selective
(b), query-driven aggregation (c), and local recency (d)—purely from gauge transport and
KL-divergence dynamics, without learned attention projections (WQ, WK , WV ).
To contextualize the gauge VFE result, we compare against the strongest purely non-
neural baseline on WikiText-103: the modified Kneser-Ney 5-gram model (Chen and Good-
man, 1998). Merity et al. (2017) report KN-5 test perplexity ∼153–156, but this uses
48
Attention and Transformers as Free Energy Minimization
word-level tokenization with a ∼267K vocabulary; since perplexity is computed per-token
and a larger vocabulary makes next-token prediction harder, this figure is not directly com-
mensurable with BPE-tokenized results. To obtain a controlled comparison, we trained a
modified Kneser-Ney 5-gram model on WikiText-103 using the same GPT-2 BPE tokeniza-
tion (50,257 vocabulary) as the gauge VFE model.1 Under matched tokenization, the KN-5
baseline achieves test perplexity 134.8 (validation 135.9), as shown in Table 6. The gauge
VFE model (test PPL 74.9) surpasses this matched baseline by 59.9 PPL (1.80×), confirm-
ing under controlled conditions that gauge-covariant variational inference outperforms the
strongest classical statistical method. The gauge VFE achieves this improvement through
an entirely different mechanism: variational inference over gauge-transported Gaussian be-
liefs rather than count-based smoothing suggesting that the geometric structure of the
framework captures statistical regularities inaccessible to n-gram models. For reference,
neural language models on WikiText-103 range from PPL ∼49 for basic LSTMs (Grave
et al., 2017) to PPL ∼18 for large optimized transformers (Dai et al., 2019), establishing
the performance ceiling for architectures with full neural components.
Closing the gap to neural baselines is not the goal of this work; rather, these results
establish that gauge-covariant variational inference is a viable computational substrate for
language modeling (one that already outperforms the best classical statistical methods)
supporting the explanatory thesis that transformers succeed because they approximate this
underlying mathematical structure.
5.3.3 Emergent Semantic Structure
We analyze the learned belief embeddings μ ∈ R90 and gauge frames ϕ ∈ L9
a=1 gl(10) ∼=
R9×100=900 for emergent categorical structure, using five linguistic categories (content word,
digit, function word, letter, punctuation) across 500 representative tokens. Clustering met-
rics are computed in the full-dimensional spaces at every 10,000 training steps.
Belief Embeddings. The learned belief means μ develop statistically significant cat-
egorical structure over training. The inter-class to intra-class distance ratio grows from
1.00 at initialization to 1.06 at convergence, indicating that tokens of different linguistic
categories become modestly more separated than tokens within the same category. Per-
dimension ANOVA tests find 82% of the 90 embedding dimensions exhibit statistically
significant category dependence (geometric mean p = 1.4 × 10−5, mean F = 7.1). The first
three principal components capture 16.9% of the variance (50% at 20 components), and
the Calinski-Harabasz index reaches 7.7. The silhouette score is essentially zero (+0.004),
indicating that the categories approach the boundary of forming separable clusters in the
full 90-dimensional space, a substantial improvement over the anti-clustering regime typical
of random embeddings.
Gauge Frames. The gauge frame coefficients ϕ ∈ R900 (9 heads × 100 dimensions per
gl(10) head) present a richer picture. PCA visualization reveals clear categorical separation
in the top principal components: punctuation tokens separate sharply from content words
1. Our matched-tokenization KN-5 implementation uses three discount parameters per order following Chen
and Goodman (1998), trained on ∼119M BPE tokens (∼212M unique n-grams across orders 1–5). The
∼18 PPL reduction from the word-level figure (∼153–156 → 134.8) is consistent with the smaller BPE
vocabulary making per-token prediction easier.
49
Dennis
along PC1, with function words, letters, and digits occupying distinct intermediate regions
(Figure ??). This low-dimensional separation is more visually striking than that of the belief
embeddings, despite the first three principal components capturing only 6.8% of the total
variance across 900 dimensions. Per-dimension ANOVA analysis confirms this structure:
62% of 100 sampled gauge frame dimensions are statistically significant (geometric mean p =
7.7 × 10−4, mean F = 4.7). Yet the full-dimensional silhouette score is negative (−0.105),
and the inter/intra class distance ratio is 1.055 (near unity), indicating that in the bulk
of the 900-dimensional space, within-category distances exceed between-category distances.
This apparent contradiction—clear low-dimensional categorical separation coexisting with
negative full-dimensional silhouette—reveals a dual structure: the gauge frames allocate a
low-dimensional subspace to encoding token category (visible in the top PCs) while devoting
the vast majority of their capacity to within-category diversification, equipping tokens of
the same linguistic type with distinct transport operators Ωij = exp(ϕi) exp(−ϕj ) suited to
their different contextual roles.
Quantitative Clustering Evidence. The clustering metrics for μ evolve progressively
throughout training, strengthening monotonically rather than peaking early and declining.
The Calinski-Harabasz index grows from 1.0 at initialization to 7.7 at convergence (peak-
ing at 7.8 at step 50,000), and the ANOVA-significant fraction increases from 1% to 82%
(peaking at 84% at step 50,000), indicating that category-driven organization strengthens
as the model optimizes for language modeling. For ϕ: the metrics stabilize after step 20,000
(silhouette ≈ −0.105 to −0.131, Calinski-Harabasz ≈ 4.0–4.6, ANOVA fraction ≈ 49–66%),
consistent with the gauge frames converging to a transport-optimized configuration and
refining within that regime.
Implications. The contrast between μ and ϕ is theoretically informative. The belief
means μ directly couple to the observation model p(o | μ) through the output projection,
creating direct pressure to encode token identity; the progressively strengthening categorical
structure (82% ANOVA significance, near-zero silhouette) reflects this coupling. The gauge
frames ϕ parameterize transport operators Ω = exp(ϕ) that mediate inter-token commu-
nication, and their dual structure—categorical separation in a low-dimensional subspace,
within-category diversification in the bulk—reflects a dual functional requirement. The
low-dimensional categorical component ensures that transport operators respect the coarse
linguistic type of a token (e.g., punctuation tokens share a qualitatively different frame
orientation from content words, visible in Figure ??). The high-dimensional diversification
ensures that tokens of the same category are equipped with distinct transport operators
suited to their specific contextual roles, which drives the negative full-dimensional silhou-
ette (−0.105). This suggests that the transport operators Ωij play a complementary role
to the attention weights βij : while the attention distribution determines which tokens com-
municate, the gauge transport determines how information is geometrically transformed
during that communication—and this “how” depends on both the categorical identity and
the contextual role of each token.
These results demonstrate that the gauge-theoretic framework scales to production lan-
guage modeling tasks, with the theoretical constructs (μ, Σ, ϕ, KL-attention) functioning
as intended during actual training.
50
Attention and Transformers as Free Energy Minimization
Representational Capacity vs. Output Capacity. The full representational capacity
is K + nheads · d2
head dimensions per token (990 for K = 90 with 9 heads of GL(10)), yet
only the K-dimensional belief means μ couple to the output. This is a structural conse-
quence of gauge covariance: the observation model must depend only on gauge-covariant
quantities, not on the internal frame parameters. The gauge frames influence predictions
indirectly through the transport operators Ωij that mediate belief aggregation during VFE
iterations. Increasing K (with multi-head GL(dhead) to control matrix exponential cost) is
the theoretically justified route to higher output capacity.
6 Discussion
6.1 Summary of Contributions
Standard attention as a degenerate limit. The BERT validation (Section 5.2) demon-
strates quantitative agreement between gauge-aligned KL attention and standard dot-
product attention across 105 diverse passages and all 144 heads (¯r = 0.804 at optimal τ ;
Figure 2), with the entropy of the two distributions closely matched (H(β)/H(α) = 1.076;
Section 5.2.4). The multi-model validation (Section 5.2.5) extends this to five architectures
(¯r > 0.6 for all, ¯r = 0.851 for ALBERT), and the sequence-length analysis (Section 5.2.6)
confirms robustness from N = 16 to N = 512 tokens. Beyond the correlation itself (which
partly reflects the algebraic relationship between squared distance and dot product un-
der approximately constant norms) the framework generates non-trivial predictions that
are confirmed: the optimal temperature matches the dimensional prediction τ ≈ 2√dk
to within 19%, the key-norm bias has the predicted sign and its temperature dependence
matches theory, and deeper layers exhibit stronger correlations consistent with increasingly
structured representations.
The unreduced framework as a language model. The GL(10) gauge VFE achieves
test perplexity 74.9 on WikiText-103 (Table 6), surpassing the matched Kneser-Ney 5-
gram baseline (PPL 134.8) by 1.80× without MLPs, pointwise activation functions, or
learned attention projections. That gauge-covariant variational inference alone produces
a functioning language model supports the explanatory claim that transformers succeed
because they approximate this underlying mathematical structure.
6.2 Architectural Implications for Standard Transformers
The framework provides theoretical explanations for empirical design choices in transform-
ers. The framework derives a key-norm cancellation condition (∥μj ∥2 ≈ C) required for
frame-independent inference (Section 4.2.1); layer normalization is one mechanism that
achieves this condition, explaining why some form of norm regulation is empirically benefi-
cial. The temperature scaling 1/√dk is the natural normalization arising from dimensional
concentration of measure. Multi-head attention corresponds to restricting the gauge algebra
to a block-diagonal subalgebra L
a gl(dhead), with each head implementing an independent
irreducible representation (Section 4.9). The output projection WO provides the only stan-
dard mechanism for post-attention cross-head information mixing.
These identifications also suggest alternatives to current practice. While layer normal-
ization mitigates key-norm bias uniformly across all tokens, the gauge-theoretic perspective
51
Dennis
suggests learning optimal token norm profiles that trade representational capacity (large
norms encode more information) against attention quality (low norms avoid bias) in a
Shannon-like manner, implemented via soft norm constraints or adaptive normalization
schedules that vary across layers or heads.
6.3 The Flat Bundle Limit and the Geometry of Language
A structural observation follows from the 0-dimensional base manifold of the transformer
architecture. The gauge transport (12),
Ωij = eϕi e−ϕj , (117)
is pure gauge by construction. To see what this implies, consider the discrete analog of
curvature: the holonomy around a closed loop of agents i → j → k → i,
Hijk = Ωij Ωjk Ωki = eϕi e−ϕj eϕj e−ϕk eϕk e−ϕi = I. (118)
The consecutive inverse–exponential pairs cancel identically: e−ϕj eϕj = I. The holonomy
vanishes for every closed loop, regardless of the learned gauge frames ϕi. The bundle is
automatically flat. This is not an approximation or a simplifying assumption rather it is a
theorem of the architecture. Any transport of the form Ωij = gi g−1
j for vertex-local group
elements gi satisfies the cocycle condition Ωij Ωjk = Ωik and therefore has vanishing discrete
curvature.
Standard transformers share this property. The learned projections WQ, WK parametrize
the gauge transport via WQW ⊤
K = σ−2Ω−⊤ (Section 2.1.3), which is position-independent,
i.e. a constant, and therefore flat, gauge transformation. Dot-product attention operates in
this flat bundle limit.
Both architectures are thus restricted to the same geometric regime: flat bundles with
trivial holonomy. In the most general discrete gauge theory on a token graph, the transport
Ωij on each edge would be an independent group element, and the holonomy around closed
loops would be unconstrained. A system with non-trivial holonomy would exhibit path-
dependent meaning transport: relating token i to token k directly via Ωik would yield a
different result than relating them via an intermediate token j through Ωij Ωjk. The flat
bundle limit eliminates this path-dependence entirely.
Conjecture (flat bundle optimality). Natural language, as an informational system,
has evolved toward minimal holonomy: that is, the semantic transport between contexts is
approximately path-independent. Transformers succeed not despite their restriction to flat
bundles, but because of it: they are architecturally matched to the geometric structure that
language has already converged to.
The intuition is that a language with non-trivial holonomy would be communicatively
pathological. If the meaning transported from speaker A to speaker C depended on which
intermediate speakers relayed the message, reliable communication would be impossible.
Selective pressure toward unambiguous, composable meaning is selective pressure toward
flat gauge structure. The cocycle condition Ωij Ωjk = Ωik is automatic in both our frame-
work and in standard transformers and furthermore is the mathematical expression of this
communicative coherence.
52
Attention and Transformers as Free Energy Minimization
This conjecture carries a testable prediction: linguistic phenomena where transformers
systematically struggle such as pragmatic inference, irony, context-dependent ambiguity,
culturally situated meaning, etc may correspond to residual curvature in the bundle of
meaning, i.e., cases where semantic transport is genuinely path-dependent and the flat bun-
dle approximation breaks down. If confirmed, measuring learned holonomy in generalized
(non-flat) gauge architectures could identify the geometric boundaries of transformer-based
language modeling.
6.4 Limitations
Our BERT validation tested the flat bundle limit where all frames are globally aligned and
the connection is trivial. This is the appropriate limit for comparison with standard trans-
formers, which lack explicit gauge structure, but it does not probe the regime where the
framework is most distinctive. While the multi-model validation (Section 5.2.5) extends the
results to five architectures (BERT-base, BERT-large, DistilBERT, RoBERTa, ALBERT),
all share dhead = 64 and English-language pretraining. Extending to autoregressive archi-
tectures (GPT, LLaMA), non-64 head dimensions, larger model scales, and non-English
languages remains future work.
The current comparison further focused on the scaled isotropic limit where uncertainty
covariances Σi = σ2I are absorbed into learned projections. The full theory incorporates
non-trivial covariance matrices Σi encoding per-token uncertainty, corresponding to “fuzzy”
vector embeddings in unique token/agent frames. Comparing this against transformers
with explicit uncertainty estimation (Bayesian neural networks, ensemble methods) remains
untested.
The gauge VFE is computationally more expensive than standard attention, primarily
due to the matrix exponential exp(ϕi · G) required for transport operators and the need for
per-pair KL divergences rather than batched dot products. We emphasize that this is an
explanatory theory: the goal is not to produce state-of-the-art language models but to pro-
vide a mathematical foundation for understanding why attention-based architectures work.
Just as statistical mechanics explains thermodynamics without replacing steam engines, the
gauge-theoretic framework explains transformers without replacing them.
6.5 Future Directions
6.5.1 Cross-Head Gauge Mixing.
Perhaps the most structurally significant generalization is the restoration of off-diagonal
gauge mixing (Section 4.4.3). Standard multi-head attention restricts the gauge algebra
to L
a gl(dhead) ⊂ gl(dk), projecting out the vast majority of generators (87.5% for typical
architectures). The projected-out generators correspond to cross-head gauge transport;
VFE gradient components ∇ϕF m that would drive the system to learn how information in
one head should be rotated into another head’s frame during the attention computation,
not merely post-attention via WO. This is qualitatively different from what the output
projection provides: off-diagonal transport makes the attention scores themselves cross-
head aware, so that head a’s decision about what to attend to can depend on the content
of head b.
53
Dennis
A practical implementation uses group-structured sparsity: only selected pairs of heads
receive off-diagonal generators, and transitively connected heads are merged into super-
blocks for the block-diagonal KL computation. This preserves most of the memory savings of
the standard block-diagonal structure (the blocks are simply larger where heads are coupled)
while enabling cross-head transport at the pairs where it matters most. For equivariant
architectures with SO(3) irrep structure, the off-diagonal blocks are precisely the Clebsch–
Gordan couplings between angular momentum channels, the physically natural mechanism
by which scalar, vector, and tensor information interact. Whether the computational cost of
richer transport can be justified by improved expressiveness or whether structured sparsity
achieves most of the benefit at reduced cost is a central experimental question.
6.5.2 Non-Uniform Head Dimensions.
Standard transformers use uniform head dimensions (dhead = d/H) rather than the geo-
metrically natural dimensions dictated by irrep structure (e.g., 2ℓ + 1 for SO(3) irreps ℓ).
Whether transformers with non-uniform head dimensions matching irrep sizes achieve bet-
ter performance or sample efficiency on tasks with known symmetries is an open question
with practical implications for equivariant architectures.
6.5.3 Gauge Symmetry as Inductive Bias.
By leveraging gauge theory coupled with information geometry, we can introduce geometric
features that constrain model behavior according to symmetry principles, analogous to how
convolutional neural networks impose translation equivariance or graph neural networks
respect permutation invariance (Finzi et al., 2020; Weiler et al., 2018; Kondor and Trivedi,
2018). The GL(K) invariance theorem allows us to interpret that standard transformers
already implement gauge symmetry compatible with KL-based attention. Our framework
shifts toward a landscape where specific subgroups (SO(N ), SU(N ), Lorentz group, etc.)
can be imposed as inductive biases, potentially enabling transformers to learn from fewer
examples in domains with known physical structure.
6.5.4 Fields of Transformers.
Standard attention transformers are a 0-dimensional gauge theory in which all geometric
structure is absorbed into learned parameters. Lifting this restriction opens a qualitatively
new construction: fields of transformers over an N -dimensional base manifold, coupled by
induced gauge connections. At each point c ∈ C, a local transformer processes tokens using
the fiber structure (beliefs, covariances, gauge frames); the gauge connection Aμ(c) defines
how information is parallel-transported between transformers at neighboring points. This
transport is generally path-dependent: the holonomy P exp− H Aμ dxμ around a closed
loop encodes curvature, meaning the path information takes between two points affects
how it is transformed upon arrival. Such fields of transformers would be natural for data
with intrinsic spatial or temporal extent (video (2D+time), multi-agent robotics (physical
space), or scientific data on manifolds) where the geometric relationships between processing
locations carry information that flat architectures discard.
54
Attention and Transformers as Free Energy Minimization
6.5.5 Multi-Timescale Dynamics and Backpropagation-Free Learning.
The full variational free energy (Section 3.4) includes both the fast belief subsystem Ffast and
the slow model subsystem Fslow, with timescale separation ηs ≪ ηq. Standard transformer
training operates entirely within this separation: the forward pass corresponds to fast-
variable dynamics, while gradient descent on weights corresponds to slow-variable learning.
Current transformers are thus the adiabatic limit where slow variables are frozen during
inference and updated only between forward passes via backpropagation.
The full VFE prescribes coupled dynamics at both timescales simultaneously. The
model-model alignment term γij DKL(si∥ ˜Ωij sj ) drives agents toward consensus not just on
what the world state is (beliefs), but on how the world works (models). Combined with hier-
archical priors that couple beliefs to model uncertainty, this yields three coupled dynamical
regimes: fast dynamics (qi) corresponding to in-context belief updating; slow dynamics (si)
corresponding to online model adaptation; and cross-scale coupling ensuring that model
uncertainty propagates into belief uncertainty. Implementing the full VFE with both sub-
systems active during inference would yield a model that treats language not as a static
distribution to be memorized but as a dynamical information system to be tracked, with
in-context learning, fine-tuning, and continual adaptation arising as limits of the same un-
derlying variational dynamics rather than requiring separate architectural mechanisms.
A natural extension toward backpropagation-free learning would update all variational
parameters (μi, Σi, ϕi) via local VFE gradient descent without backpropagation through
the computational graph. When combined with a delta rule for the output projection
Wout (∆W = η(target − prediction) ⊗ μ⊤), the entire learning pipeline would become
backpropagation-free whereby beliefs evolve by VFE minimization (fast timescale), gauge
frames via ∂F/∂ϕ (intermediate), and the output mapping via local error-driven learning
(slow). Such an approach would demonstrate whether the variational free energy principle
provides not only an explanatory framework but a viable alternative learning algorithm
grounded in information geometry.
7 Conclusion
We have presented a gauge-theoretic framework in which attention emerges as the commu-
nication mechanism of distributed variational inference. Treating language as a dynamic
informational system whereby agents exchange beliefs under uncertainty via gauge-covariant
variational free energy minimization provides the mathematically natural formalism, with
attention weights, transport operators, and gradient updates each corresponding to identi-
fiable information-theoretic quantities within a single variational objective.
The GL(K) invariance theorem (Theorem 1) reveals that standard transformers already
implement this gauge structure: the learned projections WQ, WK are gauge transformations,
and standard attention is recovered as a degenerate limit in which gauge degrees of freedom
are absorbed into learned parameters. Our empirical results validate this perspective: the
transformer validation confirms quantitative agreement between KL-based and dot-product
attention across five architectures and 105 passages (¯r = 0.804, H(β)/H(α) = 1.076; Sec-
tion 5.2), with the optimal temperature predicted from head dimension alone (τ ≈ 2√d);
and the unreduced GL(K) gauge VFE achieves non-trivial language modeling surpassing
classical baselines without neural components, with the model learning structured, head-
55
Dennis
specialized attention patterns through gauge transport and KL-divergence dynamics alone
(Section 5.3).
The framework further reveals a geometric interpretation of learning itself: the untrained
network corresponds to a gauge-symmetric vacuum state whose degeneracy is broken by
observations, with training acting as explicit symmetry breaking. We anticipate this per-
spective (that attention is gauge-covariant inference, and learning is symmetry breaking)
will find application not only in machine learning but in understanding biological cogni-
tion, collective intelligence, and other domains where distributed agents must communicate
under geometric constraints.
7.1 Code Availability
The complete simulation suite that implements the gauge-theoretic variational inference
framework described in this work is publicly available at https://github.com/cdenn016/epistemic-
geometry.
All experiments reported in the results section can be reproduced using the provided
configuration files and random number generator seeds documented in the repository. The
codebase requires Python 3.9+ with NumPy, PyTorch, SciPy, and Joblib dependencies as
specified in requirements.txt.
Acknowledgments
Claude Opus 4.5/4.6 was utilized for programming our variational free energy descent sim-
ulation suite. All code was manually reviewed, corrected, and mathematically validated by
the author by hand. Furthermore, Claude was utilized for typesetting figures, LaTeX equa-
tions, and general organizational and manuscript clarity advice. The author emphasizes
that Claude (or any other LLMs) played no role in the development or conceptual ideas of
our theoretical framework. The author further declares no funding or conflicts of interest.
Note on appendix structure. Appendix A reviews standard differential geome-
try and gauge theory to make the paper self-contained for readers from machine learning
backgrounds. Readers familiar with fiber bundles, gauge connections, and curvature may
skip directly to Appendix B (Covariance Dynamics), Appendix C (Gauge frame gradients
and numerical methods) Appendix D (Variational Gradient Descent: Implementation and
Numerical Methods), Appendix E(Conditional Uniqueness of Forward KL), Appendix F
(Symmetry breaking and 2D base manifold experiments).
Appendix A. General Mathematical Framework
A.1 Principal Bundle and Associated Bundles
The following geometric and probabilistic constructions comprise standard methods in dif-
ferential geometry and gauge theory. See references for details, proofs, and examples (Naka-
hara, 2003; Frankel, 2011; Blei et al., 2017; Amari, 2016).
Let π : N → C be a smooth principal G-bundle where C is a smooth manifold (the base
space) and G is a Lie group (the structure group) acting freely and transitively on the right
on N . The projection satisfies π(n · g) = π(n) for all g ∈ G, n ∈ N .
56
Attention and Transformers as Free Energy Minimization
Let ρq : G → Aut(Bq) and ρp : G → Aut(Bp) be representations of G on smooth sta-
tistical manifolds Bq (belief/recognition fiber) and Bp (model/prior fiber). These fibers are
typically K-dimensional probability simplices ∆K (for categorical distributions) or statisti-
cal manifolds with information-geometric structure (e.g., Gaussian manifolds, exponential
families).
The associated bundles are:
Eq := N ×ρq Bq = (N × Bq) / ∼q, (119)
Ep := N ×ρp Bp = (N × Bp) / ∼p, (120)
where (n · g, b) ∼u (n, ρu(g)b) for u ∈ {q, p}.
These give fiber bundles πEq : Eq → C and πEp : Ep → C with fibers Bq(c) ∼= Bq and
Bp(c) ∼= Bp at each c ∈ C.
A.2 Agents and Multi-Agent Systems
Agent
An agent Ai is a pair of local sections over a domain Ui ⊂ C:
Ai = σi
q, σi
p
 , (121)
where σi
q : Ui → Eq and σi
p : Ui → Ep.
We write qi(c) := σi
q(c) ∈ Bq(c) and pi(c) := σi
p(c) ∈ Bp(c) for the belief and model at
base point c.
Multi-Agent System
A multi-agent system M over C is a collection of agents indexed by I:
M = Ai = σi
q, σi
p

i∈I . (122)
Agents generally overlap on intersections Ui ∩ Uj .
Meta-Agent and Epistemic Death
A meta-agent is a multi-agent system whose component agents share identical section values
on their overlap:
qi(c) = qj (c), pi(c) = pj (c) for c ∈ Ui ∩ Uj . (123)
A set of agents is epistemically dead if they identically share both beliefs and models.
While constituent agents of a meta-agent may be epistemically dead, the meta-agent itself
need not be; such agents can be integrated out, yielding coarse-grained higher-order entities
and a route toward emergence.
A.2.1 Hierarchical Meta-Agent Emergence and Cross-Scale Coupling
The gauge-theoretic framework naturally supports the emergence of meta-agents through
coarse-graining (Shen et al., 2011). A meta-agent M(1) is a set of agents {i ∈ IM} satisfying:
57
Dennis
qi = Ωij qj ∀i, j ∈ IM (belief consensus), (124)
si = ˜Ωij sj ∀i, j ∈ IM (model consensus). (125)
When these conditions hold, the agents are epistemically dead in that they share iden-
tical beliefs and models after accounting for gauge transformations and no longer evolve
without continual observations. The meta-agent can be described by a single representative
renormalized state (qM, sM).
Meta-agents at scale ζ = 0 (individual agents) can form meta-agents at scale ζ = 1
(groups), which can further coalesce into meta-agents at scale ζ = 2 (communities), and so
on, creating a hierarchical scale structure:
Agents(0) consensus
−−−−−−→ Meta-agents(1) consensus
−−−−−−→ Meta(2) consensus
−−−−−−→ · · ·
Figure 10: Hierarchical emergence through consensus at successive scales ζ.
At each scale, the effective free energy takes the same form as Eq. (31), with agents at
scale ζ replaced by meta-agents at scale ζ + 1, coupling constants renormalized according
to
β(ζ+1)
ij = fβ ({β(ζ)
kl }), γ(ζ+1)
ij = fγ ({γ(ζ)
kl }),
and effective gauge frames given by ϕ(ζ+1)
M = average({ϕ(ζ)
i }). This is the gauge-theoretic
analogue of renormalization group flow in statistical field theory (Anderson and Rosenfeld,
1988; Wilson and Kogut, 1974; Garc´ıa-Mill´an and Pruessner, 2024). Cross-scale couplings
Λζ
ζ′ mediate interactions between agents at different hierarchical levels via the cross-scale
transport operators defined below.
Whether this hierarchical structure maps onto multi-layer transformer architectures
(where each layer constitutes a “scale”), or onto distinct organizational levels within a
single attention layer, remains an open question for future investigation.
A.3 Bundle Morphisms and Transport Operators
Via standard horizontal lifting from the principal bundle N to the associated bundles, we
obtain a hierarchy of morphisms and induced connections:
Intra-bundle transport operators Ω(q)
ij : Γ(Bq) → Γ(Bq) (belief-to-belief) and Ω(p)
ij :
Γ(Bp) → Γ(Bp) (model-to-model) act within a single bundle. Cross-scale transport opera-
tors Λs
s′ : Γs(Bq) → Γs′
(Bq) (beliefs across scales) and ˜Λs
s′ : Γs(Bp) → Γs′
(Bp) (models across
scales) connect different hierarchical levels. Inter-bundle morphisms Θi
j : Γ(Bq) → Γ(Bp)
(belief to model) and ˜Θi
j : Γ(Bp) → Γ(Bq) (model to belief) relate the two fiber types. Fi-
nally, global bundle morphisms Φ : Ep → Eq (model bundle to belief bundle) and ˜Φ : Eq → Ep
(belief bundle to model bundle) act on the total associated bundles.
Here Γ(Bu) denotes the space of smooth sections over the base C.
58
Attention and Transformers as Free Energy Minimization
A.4 Gauge Frames and Connections
Each agent i possesses a local gauge frame field:
ϕi : Ui → g = Lie(G), (126)
which induces a local connection one-form:
A(i)
μ (c) = U −1
i (c)∂μUi(c), (127)
where Ui(c) = exp[ϕi(c)] ∈ G. The associated field strength (gauge curvature) is:
F (i)
μν (c) = ∂μA(i)
ν − ∂ν A(i)
μ + [A(i)
μ , A(i)
ν ] ∈ g. (128)
When two agents overlap at c ∈ Ui ∩ Uj , the inter-agent gauge transformation:
Ωij (c) = exp[ϕi(c)] exp[−ϕj (c)] ∈ G (129)
transports agent j’s representations into agent i’s frame:
qj (c) 7 → Ωij (c) · qj (c) := ρ(Ωij (c))qj (c). (130)
On overlaps, local connections are related by:
A(i)
μ = Ωij A(j)
μ Ω−1
ij + Ωij ∂μΩ−1
ij . (131)
Appendix B. Covariance Dynamics and Equilibrium Analysis
B.1 Covariance Gradient of the Generalized Free Energy
Here we derive the gradient of the single-agent free energy Fi with respect to the covariance
Σi for agent i; a well known result in information geometry.
The free energy decomposes as
Fi = DKL(qi ∥ pi) + X
j̸ =i
βij DKL(qi ∥ Ωij qj ) − Eqi [log p(oi | ki)], (132)
where qi = N (μi, Σi), pi = N (μp,i, Σp,i), and P
j βij = 1 by construction.
B.1.1 Gaussian KL Divergence under GL(K) Gauge Transport
The alignment term DKL(qi ∥ Ωij qj ) involves KL divergences between Gaussians related by
gauge transport Ωij = eϕi e−ϕj ∈ GL+(K). The standard Gaussian KL divergence is
DKL(N (μ1, Σ1) ∥ N (μ2, Σ2)) = 1
2
h
log |Σ2|
|Σ1| + trΣ−1
2 Σ1
 (133)
+ (μ2 − μ1)⊤Σ−1
2 (μ2 − μ1) − d
i
. (134)
and differentiating w.r.t. Σ1 (holding μ1, μ2, Σ2 fixed) gives
59
Dennis
∂DKL
∂Σ1
= 1
2
−Σ−1
1 + Σ−1
2
 . (135)
Applying this to each term in Fi, where the alignment KL involves the GL(K)-transported
belief Ωij qj = N (Ωij μj , Ωij Σj Ω⊤
ij ):
∂
∂Σi
DKL(qi ∥ pi) = 1
2
h
−Σ−1
i + Σ−1
p,i
i
, (136)
∂
∂Σi
DKL(qi ∥ Ωij qj ) = 1
2
h
−Σ−1
i + (Ωij Σj Ω⊤
ij )−1i
. (137)
The observation term contributes an O(Σ−1
i ) correction that we neglect in the high-
precision / strong-alignment regime. The alignment term P
j βij DKL(qi ∥ Ωij qj ) requires
the product rule, since βij = softmaxj (−DKL/τ ) depends on Σi through the KL divergence.
Thus
∂Fi
∂Σi
= 1
2

−2Σ−1
i + Σ−1
p,i + X
j
βij (Ωij Σj Ω⊤
ij )−1


+ X
j
∂βij
∂Σi
DKL(qi ∥ Ωij qj ).
(138)
The first line contains the “direct” gradient terms (differentiating each Kij at fixed βij ).
The second line is the attention weight gradient, analogous to the ∂βij /∂μi term in the
mean gradient. From the softmax structure, P
j ∂βij /∂Σi = 0, so this correction vanishes
when the KL divergences are uniform across neighbors, and also in both the τ → 0 and
τ → ∞ limits.
Because P
j βij = 1, there is no (1 + P
j βij ) prefactor in front of Σ−1
i .
B.2 Fixed-Point Equation and Symmetric Solution
At equilibrium, we set ∂Fi/∂Σi = 0. Since P
j ∂βij /∂Σi = 0, the attention weight cor-
rection term vanishes whenever the KL divergences DKL(qi ∥ Ωij qj ) are approximately uni-
form across neighbors, and also in both the hard-attention (τ → 0) and uniform-attention
(τ → ∞) limits. In these regimes, the equilibrium condition reduces to
Σ−1
i = 1
2

Σ−1
p,i + X
j
βij (Ωij Σj Ω⊤
ij )−1

 . (139)
This is a matrix-valued fixed-point equation coupling all agents: each agent’s preci-
sion is the β-weighted combination of its own prior precision and the transported neighbor
precisions. At intermediate temperatures, the attention weight correction provides an addi-
tional contribution that couples precision dynamics to the variance of KL divergences across
neighbors.
60
Attention and Transformers as Free Energy Minimization
B.2.1 Homogeneous limit
In the homogenous limit we assume
(i) all agents are identical, so Σi = Σ∞ for all i,
(ii) Ωij ≈ I (weak misalignment),
(iii) Σp,i = Σ0 (shared prior).
Then (139) becomes
Σ−1
∞ = 1
2
Σ−1
0 + Σ−1
∞
 =⇒ Σ∞ = Σ0. (140)
Hence, in a perfectly symmetric population, the equilibrium covariance reproduces the
shared prior.
B.2.2 Alignment-dominated regime
Although P
j βij = 1, the effective strength of alignment is controlled by the parameter τ
in
βij = exp− 1
τ DKL
qi ∥ Ωij qj

P
k exp− 1
τ DKL
qi ∥ Ωikqk
 . (141)
As τ → 0, βij becomes sharply peaked on whichever neighbor j best agrees (after
transport). In that low-τ limit, the prior precision Σ−1
p,i becomes negligible relative to the
socially enforced alignment term, and
Σ−1
i ≈ X
j
βij (Ωij Σj Ω⊤
ij )−1 = (Ωij Σj Ω⊤
ij )−1
β . (142)
Here ⟨·⟩β denotes a β-weighted expectation over neighbors.
Therefore, in the strong-alignment (i.e. small τ ) regime, agent i’s precision matrix
becomes the β-weighted average of its neighbors’ transported precisions.
When, additionally, all agents already have approximately equal covariances and the
transports are near-identity, i.e. Σi ≈ Σj and Ωij ≈ I, then (142) implies
Σ−1
i ≈ Σ−1
j =⇒ Σi ≈ Ωij Σj Ω⊤
ij , (143)
justifying the alignment assumption used in the main text.
B.2.3 Gradient flow dynamics
Finally, consider the gradient flow
dΣi
dt = −ηΣ
∂Fi
∂Σi
, ηΣ > 0. (144)
Local stability of the equilibrium (139) follows from the positive-definiteness of the
Hessian. For the Gaussian KL terms,
∂2DKL
∂Σ1 ∂Σ1
∼ Σ−1
1 ⊗ Σ−1
1 + Σ−1
2 ⊗ Σ−1
2 , (145)
61
Dennis
which is manifestly positive definite for Σ1, Σ2 ≻ 0.
Hence the covariance alignment fixed-point is an attractor of the variational dynamics.
Therefore, we find that Σi ≈ Ωij Σj Ω⊤
ij emerges from the dynamics itself rather than as
an imposed constraint.
Appendix C. Gauge Frame Gradients
We derive the gradient of the variational free energy F with respect to the gauge frame
parameters ϕi ∈ g. Unlike the covariance gradients (Appendix B), gauge frame gradients
require the differential of the matrix exponential map, and the resulting expressions differ
for compact (e.g., SO(N )) and general (GL(K)) gauge groups.
C.1 Structure of the Gauge Frame Gradient
The gauge frames enter the free energy exclusively through the parallel transport operators
Ωij = exp(ϕi) exp(−ϕj ). Since neither the belief-prior term DKL(qi∥pi) nor the observation
likelihood −Eqi [log p(oi|ki)] depends on ϕi, the complete gradient is:
∂F
∂ϕi
= X
j
 ∂βij
∂ϕi
Kij + βij
∂Kij
∂ϕi

+ X
k
"
βki
∂Kki
∂ϕi
+ X
l
∂βkl
∂ϕi
Kkl
#
, (146)
where Kij ≡ DKL(qi∥Ωij qj ). The first sum captures agent i aligning to others, with
the product rule applied to βij Kij . The second sum captures others aligning to i: chang-
ing ϕi shifts Kki through Ωki = exp(ϕk) exp(−ϕi), which both directly adjusts the align-
ment cost (the βki ∂Kki/∂ϕi term) and redistributes all of agent k’s attention weights (theP
l(∂βkl/∂ϕi) Kkl term). Since only Kki depends on ϕi in agent k’s softmax, the latter
simplifies to:
X
l
∂βkl
∂ϕi
Kkl = − βki
τ
∂Kki
∂ϕi
Kki − X
l
βklKkl
!
. (147)
This correction vanishes when Kki equals agent k’s average divergence ¯Kk = P
l βklKkl,
or when βki → 0 (agent k does not attend to i).
The coupling weight gradient in the first sum follows from the softmax structure:
∂βij
∂ϕi
= − βij
τ
"
∂Kij
∂ϕi
− X
k
βik
∂Kik
∂ϕi
#
. (148)
C.2 Differential of the Matrix Exponential
All KL gradient terms require the derivative of the transport operator with respect to the
gauge frame. For Ωij = exp(ϕi) exp(−ϕj ):
62
Attention and Transformers as Free Energy Minimization
∂Ωij
∂ϕi
= dexpϕi (·) · exp(−ϕj ), (149)
∂Ωij
∂ϕj
= − exp(ϕi) · dexpϕj (·) · exp(−ϕj ), (150)
where dexpϕ : g → g denotes the differential of the exponential map at ϕ. For a direction
ξ ∈ g, the differential is given by (??):
dexpϕ(ξ) =
∞X
n=0
1
(n + 1)! adn
ϕ(ξ) = eadϕ − I
adϕ
(ξ), (151)
where adϕ(ξ) = [ϕ, ξ] = ϕξ − ξϕ is the adjoint representation.
C.2.1 SO(N ) specialization.
For SO(N ), the gauge frame is expanded in skew-symmetric generators: ϕ = P
a ϕaTa
where Ta ∈ so(N ), and the gradient is computed component-wise as ∂F/∂ϕa. For SO(3)
with ∥ϕ∥ = θ, the differential admits a closed form via the Rodrigues formula:
dexpϕ(Ta) = Ta + c1(θ)[ϕ, Ta] + c2(θ)[ϕ, [ϕ, Ta]], (152)
with coefficients:
c1(θ) = 1 − cos θ
θ2 , c2(θ) = θ − sin θ
θ3 . (153)
For θ < ϵ (numerically ϵ ∼ 10−4), Taylor expansions c1 ≈ 1
2 − θ2
24 and c2 ≈ 1
6 − θ2
120
provide stable evaluation.
C.2.2 GL(K) specialization.
For GL(K), the gauge frame ϕ ∈ gl(K) ∼= RK×K is a general K ×K matrix. The differential
of the matrix exponential does not simplify to a Rodrigues-like formula, and is instead
computed via the integral representation:
dexpϕ(ξ) =
Z 1
0
etϕ ξ e(1−t)ϕ dt, (154)
or equivalently through the (2K × 2K) block matrix identity (?):
exp
ϕ ξ
0 ϕ

=
eϕ dexpϕ(ξ)
0 eϕ

. (155)
In practice, automatic differentiation through torch.matrix exp provides the GL(K)
differential without explicit computation of (154) or (155).
63
Dennis
C.3 KL Gradient Through Transport
The KL divergence Kij = DKL(qi∥Ωij qj ) between qi = N (μi, Σi) and the transported belief
Ωij qj = N (Ωij μj , Ωij Σj Ω⊤
ij ) depends on ϕi through Ωij . Using the Gaussian KL (134) and
the chain rule:
∂Kij
∂ϕa
i
= ∂
∂ϕa
i
(
1
2
"
log |Ωij Σj Ω⊤
ij |
|Σi| + tr

(Ωij Σj Ω⊤
ij )−1Σi

+ ∆μ⊤
ij (Ωij Σj Ω⊤
ij )−1∆μij − K
#)
,
(156)
where ∆μij = μi − Ωij μj . Let ˜Σij ≡ Ωij Σj Ω⊤
ij and ˜μij ≡ Ωij μj denote the transported
statistics. The derivative has three contributions:
Mean term. From the Mahalanobis distance:
∂
∂ϕa
i
h
(μi − ˜μij )⊤ ˜Σ−1
ij (μi − ˜μij )
i
= −2(μi − ˜μij )⊤ ˜Σ−1
ij
∂ ˜μij
∂ϕa
i
+(terms from ∂ ˜Σ−1
ij /∂ϕa
i ). (157)
Trace term. From the precision-covariance product:
∂
∂ϕa
i
tr( ˜Σ−1
ij Σi) = −tr ˜Σ−1
ij
∂ ˜Σij
∂ϕa
i
˜Σ−1
ij Σi
!
. (158)
Log-determinant term. From the log-determinant ratio:
∂
∂ϕa
i
log | ˜Σij | = tr ˜Σ−1
ij
∂ ˜Σij
∂ϕa
i
!
. (159)
Transport derivatives. The derivatives of the transported statistics with respect to ϕa
i
are:
∂ ˜μij
∂ϕa
i
= ∂(Ωij μj )
∂ϕa
i
= dexpϕi (Ta) · exp(−ϕj ) μj ≡ Q(i)
a exp(−ϕj ) μj , (160)
∂ ˜Σij
∂ϕa
i
= Q(i)
a exp(−ϕj )Σj exp(−ϕj )⊤Ω⊤
ij + Ωij exp(−ϕj )Σj exp(−ϕj )⊤(Q(i)
a )⊤, (161)
where Q(i)
a ≡ dexpϕi (Ta) is the differential of the exponential map evaluated at ϕi in the
direction of generator Ta.
C.4 Retraction and Numerical Considerations
Lie algebra updates. Since ϕi ∈ g (a vector space), gradient descent proceeds in the Lie
algebra directly:
ϕi ← ϕi − ηϕ
∂F
∂ϕi
, (162)
64
Attention and Transformers as Free Energy Minimization
where ηϕ is the gauge frame learning rate. No metric correction (Fisher-Rao or other-
wise) is applied: the exponential map exp : g → G provides natural coordinates for the Lie
group, and the gradient naturally lives in g.
SO(N ) retraction. For compact groups such as SO(N ), the exponential map is periodic.
To ensure unique parameterization and avoid branch cuts, we retract to the principal ball
∥ϕ∥ < π − ϵ (with ϵ ∼ 10−2) after each update. If ∥ϕ∥ ≥ π, the angle is wrapped modulo
2π and, if necessary, the axis is flipped via antipodal identification (ϕ → −ϕ when θ > π),
exploiting the fact that exp(ϕ) = exp(ϕ − 2πˆn) for SO(3).
GL(K) updates. For GL(K), the matrix exponential is surjective onto GL+(K) (the
identity component), and no retraction is required beyond standard gradient clipping for
numerical stability. In our implementation, automatic differentiation through PyTorch’s
matrix exp handles the GL(K) differential implicitly.
C.5 Gauge Frame Preconditioning for GL(K)
While the basic update (162) uses the Euclidean gradient in the Lie algebra, the non-
compact structure of GL(K) introduces a numerical asymmetry that motivates gauge-aware
preconditioning. The Lie algebra decomposes via the Cartan involution θ(X) = −X⊤ as:
gl(K) = so(K)
| {z }
compact (skew-symmetric)
⊕ Sym(K)
| {z }
non-compact (symmetric)
⊕ R|{z}
center (trace)
, (163)
where so(K) = {X : X + X⊤ = 0} generates rotations and Sym(K) = {X : X =
X⊤, tr X = 0} generates volume-preserving shears. The critical observation is that these
two subspaces have qualitatively different gradient geometry: for the compact directions,
∥d expϕ(Ta)∥ = O(1) uniformly in ∥ϕ∥, while for the symmetric directions, ∥d expϕ(Ta)∥ ∼
exp(∥ϕ∥). Euclidean gradient descent therefore systematically over-steps in the non-compact
directions relative to the compact ones, leading to numerical instability for GL(K) gauge
frames.
We implement four preconditioning strategies of increasing geometric fidelity:
1. Norm clipping (baseline). Simple gradient norm clipping ˜g = g · min(1, c/∥g∥) with
threshold c. No geometric awareness; treats all Lie algebra directions uniformly.
2. Cartan preconditioning. Decompose the gradient into compact and non-compact
components and selectively dampen the latter. Given Lie algebra generators {Ta}ngen
a=1 with
Gram matrix Gab = tr(T ⊤
a Tb) and symmetric inner product Sab = tr(TaTb), the projector
onto the symmetric subspace is:
Psym = 1
2 G−1(G + S), (164)
and the preconditioner is M = I − (1 − λsym)Psym with dampening factor λsym ∈ (0, 1)
(typically 0.1, corresponding to 10× dampening of symmetric directions). The precondi-
tioned gradient is ˜∇ϕF = M · ∇ϕF.
65
Dennis
3. Killing form natural gradient. The Killing form of gl(K) is κ(X, Y ) = 2K tr(XY )−
2 tr(X) tr(Y ). In the generator basis, this yields the metric:
˜gab = 2K tr(T ⊤
a Tb) − 2 tr(Ta) tr(Tb), (165)
which is positive definite on sl(K) = ker(tr) and degenerate on the center R · I (regu-
larized by ˜g → ˜g + ϵI in practice). The natural gradient ˜g−1∇ϕF has no free parameters:
the metric is determined entirely by the algebra structure. This automatically handles
the compact/non-compact asymmetry since the Killing form assigns different “lengths” to
rotations versus shears.
4. Pullback natural gradient (position-dependent). The most geometrically prin-
cipled approach uses the pullback of the Gram metric through the exponential map. The
differential of exp at ϕ ∈ g in direction Ta is:
d expϕ(Ta) = exp(ϕ) · Ψ(adϕ)(Ta), Ψ(z) = ez − 1
z =
∞X
k=0
zk
(k + 1)! , (166)
where adϕ(X) = [ϕ, X] is the adjoint action and Ψ is the inverse Bernoulli function.
The pullback metric on g is:
Gab(ϕ) = ⟨Ψ(adϕ)(Ta), Ψ(adϕ)(Tb)⟩G , (167)
where ⟨X, Y ⟩G = tr(X⊤Y ) is the Frobenius inner product. At ϕ = 0, Gab(0) = Gab (the
Gram matrix). For large ∥ϕ∥ in symmetric directions, G grows exponentially, automatically
compensating the exponential amplification of d exp in the non-compact sector. The nat-
ural gradient G(ϕ)−1∇ϕF is position-dependent and exact (up to series truncation of Ψ),
requiring the structure constants f c
ab defined by [Ta, Tb] = P
c f c
abTc.
The computational cost of each mode scales as: clipping O(ngen), Cartan O(n2
gen),
Killing O(n2
gen), pullback O(n3
gen) (dominated by the linear solve G−1g). For the GL(10)
model (ngen = 100), Killing form preconditioning adds negligible overhead. The pullback
metric becomes advantageous when gauge frames develop large norms during training, as the
position-dependent compensation prevents the gradient from overshooting in the symmetric
directions where exp(ϕ) amplifies exponentially.
Appendix D. Variational Gradient Descent: Implementation and
Numerical Methods
D.1 Natural Gradient Descent on the Gaussian Manifold
A critical aspect of our implementation is the use of natural gradient descent for the sta-
tistical parameters, which exploits the information-geometric structure of the Gaussian
manifold. Rather than treating the covariance matrices Σi as Euclidean variables, we rec-
ognize them as points on the manifold of symmetric positive-definite matrices, equipped
with the Fisher-Rao metric. The natural gradient at a point Σ on this manifold is obtained
by applying the inverse Fisher metric to the Euclidean gradient, yielding an intrinsic tan-
gent vector that respects the manifold’s curved geometry. Given Euclidean gradients ∇μ
66
Attention and Transformers as Free Energy Minimization
and ∇Σ, the natural gradient projections are
δμ = −Σ ∇μ, δΣ = −2 Σ sym(∇Σ) Σ, (168)
where sym(M ) = 1
2 (M + M ⊤).
The covariance update δΣ is then passed to a manifold retraction (Section D.2) that
whitens it into the eigenbasis of Σ:
B = Λ−1/2 (V ⊤δΣ V ) Λ−1/2, (169)
where Σ = V ΛV ⊤ is the eigendecomposition. This whitening ensures that the resulting step
is affine-invariant—the same step size produces comparable effects regardless of the current
scale of the covariance matrix. Following the computation of the whitened tangent B, we
apply a trust-region constraint by clipping its Frobenius norm to a maximum relative step
size ρ. This prevents excessively large updates that could destabilize the manifold structure
or lead to ill-conditioned covariances.
D.2 Manifold Retraction for Covariance Matrices
The manifold retraction is the process of mapping the tangent vector back to a valid point on
the manifold. This is performed using the affine-invariant exponential map. Working in the
eigenbasis of Σk with eigenvalues λj , the whitened tangent B is diagonalized B = U ΛB U ⊤
and the retraction is
Σk+1 = V Λ1/2 U exp(τ ΛB ) U ⊤ Λ1/2 V ⊤, (170)
where τ is the step size and the component-wise exponentials exp(τ λ(j)
B ) are clipped to
|τ λ(j)
B | ≤ 50 to prevent floating-point overflow. This procedure is provably SPD-preserving
provided Σk is SPD and B is symmetric, which our implementation guarantees through
explicit symmetrization at multiple stages.
Additional safeguards include post-retraction sanitization via sanitize sigma, which
symmetrizes the result, raises on any floating-point anomalies (NaNs), and applies a spec-
tral floor to the eigenvalues to ensure they remain above a minimum threshold ϵSPD = 10−4,
together with a condition-number cap κmax = 104. This spectral regularization is essen-
tial for preventing numerical collapse of the covariance matrices during prolonged gradient
descent, particularly in regions where the free-energy landscape becomes very flat.
D.3 Gauge Frame Dynamics on GL(K)
For the gauge frame fields ϕi, which live in the Lie algebra gl(K) (or its subalgebra so(3) in
the fundamental representation), a different geometric structure must be respected. These
fields parameterize the local gauge frames as Ui = exp(P
a ϕa
i Ta), where {Ta} are the
Lie algebra generators, and the transport operators take the form Ωij = exp(ϕi) exp(−ϕj ).
Gradients of the alignment energy βij DKL(qi∥Ωij [qj ]) with respect to ϕi are obtained via the
chain rule through the transport operator: the KL gradient with respect to Ω is contracted
against the differential ∂Ωij /∂ϕa
i .
The differential of the matrix exponential relates infinitesimal variations δϕ to variations
in the group element via
∂
∂ϕa exp
P
b ϕbTb

= exp(X) · Qa, (171)
67
Dennis
where X = P
b ϕbTb and Qa = d expX (Ta) is the derivative of the exponential map applied
to the generator Ta. For the K = 3 fundamental representation of so(3), we use the exact
closed-form expression
Qa = Ta − c1(θ) adX (Ta) + c2(θ) ad2
X (Ta), (172)
where θ = ∥ϕ∥, c1(θ) = (1 − cos θ)/θ2, and c2(θ) = (θ − sin θ)/θ3, with Taylor-expanded
small-angle limits for θ < 10−4. For higher-dimensional representations (K > 3), where
the Cayley–Hamilton truncation requires K terms, we instead compute Qa via the Fr´echet
derivative of the matrix exponential using symmetric quadrature.
The resulting Euclidean gradient ∂F/∂ϕa can optionally be preconditioned using one
of several Lie-algebraic metrics. In the transformer implementation, these include: (i) the
Killing form natural gradient, which uses the Cartan-involution-modified Killing form ˜gab =
2K tr(T ⊤
a Tb)−2 tr(Ta)tr(Tb) as a position-independent metric with no free parameters; (ii) a
Cartan decomposition preconditioner that dampens the non-compact (symmetric) directions
of gl(K) relative to the compact so(K) directions; and (iii) the full pullback natural gradient
Gab(ϕ) = ⟨Ψ(adX )(Ta), Ψ(adX )(Tb)⟩, where Ψ(z) = (ez −1)/z, which captures the position-
dependent curvature of the exponential map.
Once the gradient direction is determined, we apply a step with learning rate ηϕ (typ-
ically 10−1), forming a candidate update ϕcand = ϕ + ηϕ δϕ. However, the Lie algebra
has a natural periodic boundary—rotation angles wrap at ±π—and numerical drift can
cause ϕ to wander outside the principal domain. To prevent this, we apply a retraction via
retract to principal ball, which wraps angles modulo 2π, applies an antipodal flip for
angles exceeding π, and clamps the result to ∥ϕ∥ < π − εmargin with εmargin = 10−2 to avoid
numerical instability near the critical points where the exponential map degenerates.
Gradient validation. All gauge frame gradient implementations are validated against
finite-difference approximations (∂F/∂ϕa
i ≈ [F(ϕi + ϵea) − F(ϕi − ϵea)]/(2ϵ)) with relative
error < 10−5.
Appendix E. Conditional Uniqueness of the Forward KL Divergence via
Variational Duality
Intuition. Why must the divergence measuring inter-agent communication cost be the
forward KL, and not some other statistical distance? The answer emerges from the in-
terplay of three requirements that any well-behaved multi-agent variational system must
satisfy simultaneously. First, we want each agent’s updated belief q∗
i to remain a tractable
exponential-family distribution (i.e. a Gaussian), so that the belief update has a closed-form
Boltzmann/Gibbs structure—a geometric mean of the prior and the transported neighbor
beliefs. Second, the attention weights βij should act as proper dual variables: the marginal
cost of attending to neighbor j should equal the divergence between the agent’s belief and
that neighbor’s signal, consistent with the envelope theorem. Third, the divergence should
be local and enter the free energy linearly in βij , preserving the softmax form of attention.
These three conditions, each individually nonrestrictive, are jointly restrictive. Alter-
native divergences violate at least one of the conditions. For example, the reverse KL
introduces 1/qi terms that break exponential-family closure, forcing transcendental equa-
tions solvable only via the Lambert W function. The χ2-divergence produces algebraic
68
Attention and Transformers as Free Energy Minimization
(polynomial-ratio) rather than exponential solutions. Symmetric divergences like Jensen–
Shannon mix logarithmic and reciprocal terms, destroying log-linearity. Only the forward
KL (with generator f (t) = t log t−t+1 yields f ′(t) = log t) produces a stationarity condition
that is linear in log qi, which is the condition for the solution to remain in the exponential
family. The following theorem makes this precise.
Within a broad, but well-defined, class, the forward KL divergence DKL(qi ∥ Ωij qj ) is
the only divergence that yields a closed-form Gibbs-type solution for the belief update and
a consistent dual interpretation for the attention weights. Agents locally minimize their
agent-specific variational free energy and the system of agents minimize their collective
global free energies. This local-global coordination gives rise to the expected forward KL
attention term.
In what follows we restrict to our ”matched-fiber” framework. The full generalization
follows by applying the appropriate bundle morphisms/intertwiners.
The uniqueness of the term is conditional and follows from three assumptions: (i) D is
local in c, of f-divergence form
Z
qi(c) f
 qi(c)
Ωij (c)qj (c)

dc; (ii) the coupling is linear; and
(iii) the minimizing belief q∗
i remains in the exponential-family (log-linear).
E.0.1 The Coupled Variational Problem
Each agent i minimizes a local free-energy functional:
Fi[βi] = min
qi
n
DKL(qi ∥ pi) + X
j̸ =i
βij D(qi, qj )
o
, (173)
For convenience, we can decompose the KL as:
DKL(qi ∥ pi) = ⟨Hi⟩qi + S(qi) + const, (174)
where
Hi(c) := − log pi(c) (local ”energy”), (175)
S(qi) := −
Z
qi(c) log qi(c) dc (Shannon entropy), (176)
⟨Hi⟩qi :=
Z
qi(c)Hi(c) dc (expected energy). (177)
This yields DKL(qi ∥ pi) = R qi(c)[log qi(c) − log pi(c)] dc.
The attention weights βij ≥ 0 (with P
j βij = 1) then are chosen by optimizing
Ji(βi) = X
j̸ =i
βij Cij + τ X
j̸ =i
βij log βij , (178)
where Cij denotes the marginal cost of attending to agent j as we’ve previously de-
scribed.
Later we will identify this as
69
Dennis
Cij := ∂Si
∂βij
, (179)
such that attention weights allocate in proportion to marginal coordination penalties.
E.0.2 The Forward KL and the Geometric-Mean Solution
Let D(qi, qj ) = DKL(qi ∥ Ωij qj ), where
δDKL(qi ∥ Ωij qj )
δqi(c) = log qi(c)
Ωij (c)qj (c) + 1.
The stationarity condition for (173) is
Hi(c) + log qi(c) + X
j
βij

log qi(c)
Ωij (c)qj (c) + 1

= λi, (180)
with λi enforcing normalization.
Solving for qi(c) then yields a Boltzmann distribution whose mean field is the geometric
average of the transported neighbor beliefs.
q∗
i (c) = 1
Zi
e−Hi(c)/2 Y
j
[Ωij (c)qj (c)]βij /2, (181)
This structure is preserved only for the forward KL divergence; alternative divergences
destroy the log-linearity of the exponent as previously discussed.
E.0.3 Dual Relation via the Envelope Theorem
At the stationary value q∗
i , the envelope theorem implies
∂Fi
∂βij
= D(q∗
i , qj ) = DKL(q∗
i ∥ Ωij qj ), (182)
such that the cost of increasing attention to j equals the forward KL divergence between
the updated belief q∗
i and the transported neighbor Ωij qj .
This then identifies the attention cost with the KL.
E.0.4 Conditional Uniqueness Theorem
Theorem. Let D(q, p) be a convex f -divergence of the form
D(q, p) =
Z
p(c) f
 q(c)
p(c)

dc,
with f convex, f (1) = 0, and, without loss of generality2, f ′(1) = 0. Suppose D is linear
inside the free energy (173) and that the attention weights satisfy P
j βij = 1. Then the
stationary solution q∗
i assumes the geometric-mean Boltzmann form
2. Any f -divergence is invariant under f (t) 7 → f (t) + c(t − 1) when q and p are normalized, so we may
always choose the representative with f ′(1) = 0.
70
Attention and Transformers as Free Energy Minimization
q∗
i (c) ∝ pi(c)1/2 Y
j
Ωij (c) qj (c)βij /2 (183)
for all choices of prior pi and neighbor beliefs {qj }, if, and only if, D is the forward
Kullback–Leibler divergence (up to the affine equivalence above):
D(q, p) = DKL(q ∥ p).
Proof. Step 1. The variational derivative of a general f -divergence.
With D(qi, Ωij qj ) = R (Ωij qj ) fqi/(Ωij qj ) dc and ratio r(c) = qi(c)/Ωij (c) qj (c), the
functional derivative is
δD
δqi(c) = f ′(r(c)). (184)
The stationarity condition δFi/δqi(c) = λ (with Lagrange multiplier λ) is then
log qi(c)
pi(c) + 1 + X
j
βij f ′
 qi(c)
Ωij (c) qj (c)

= λ. (185)
Step 2. Forward implication (DKL ⇒ geometric mean).
Next, for D = DKL, the generator is f (t) = t log t − t + 1, yielding f ′(t) = log t.
Substituting into (185) we have:
log qi − log pi + 1 + X
j
βij
log qi − log Ωij qj
 = λ.
However, since P
j βij = 1, the log qi terms become 2 log qi:
2 log qi = log pi + X
j
βij log(Ωij qj ) + (λ − 1).
By exponentiating we find the geometric-mean form (183).
Step 3. Reverse implication (geometric mean ⇒ DKL).
Assume the minimizer has the geometric mean form (183) for all priors pi and neighbor
beliefs {qj }.
Taking logarithms yields:
2 log q∗
i (c) = log pi(c) + X
j
βij logΩij (c) qj (c) + C. (186)
Next, rearranging and using P
j βij = 1 we have:
log q∗
i
pi
+ X
j
βij log q∗
i
Ωij qj
= C′. (187)
Comparing (187) with the general stationarity condition (185) term-by-term, we then
require
71
Dennis
f ′
 qi
Ωij qj

= log qi
Ωij qj
+ k (188)
for some constant k, and for all values of the ratio t = qi(c)/[Ωij (c) qj (c)]. Since the
geometric-mean form holds for all choices of pi and {qj }, the ratio t ranges over all of R+,
therefore (188) must hold globally: f ′(t) = log t + k for all t > 0.
Next, by integrating we have:
f (t) = t log t − t + kt + C0.
If we impose f ′(1) = 0 then k = 0. Similarly by imposing f (1) = 0 then C0 = 1.
Therefore, we have
f (t) = t log t − t + 1,
which is the generator of DKL(q ∥ p). □
E.0.5 Dual Cost Identification
Having established DKL is the unique divergence, the envelope theorem applied to the βij
optimization confirms that the marginal attention cost is
Cij ≡ ∂Fi
∂βij q∗
i
= DKL(q∗
i ∥ Ωij qj ), (189)
identifying the cost of attention with the information divergence between the agent’s
updated belief and the neighbor’s transported signal.
E.0.6 Interpretations
From the perspective of gauge invariance, the comparison is made between qi and the
transported Ωij qj within the same frame (and similarly for j → i); gauge covariance fixes
what is compared, while variational duality fixes how it is compared. From the perspective
of this duality, the forward KL is the only divergence that simultaneously yields a closed-
form Boltzmann solution for qi and a consistent dual cost Cij = ∂Fi/∂βij . From the
perspective of information geometry, the forward KL is the Bregman divergence generated
by the negative entropy potential, whose Hessian induces the Fisher-Rao metric and yields
the mixed/exponential family-projection Pythagorean theorem. These global properties are
not shared by generic f-divergences.
E.0.7 Summary
Under the natural assumptions of locality, linear coupling, and exponential-family closure,
the forward KL divergence
Cij = DKL(qi ∥ Ωij qj )
is a necessary consequence of the variational and geometric structure underlying agent-
agent coordination. By Theorem 1, the KL divergence is invariant under the full general
72
Attention and Transformers as Free Energy Minimization
Figure 11: Symmetric vacuum state: μcenter
Q tracking summary without observations. Top:
Norm trajectories for all six agents, converging to a shared equilibrium value. Middle left:
Symmetry measure Var(∥μ∥) decays to machine precision, confirming gauge symmetry is
preserved. Middle right: Population-averaged norm stabilizes. Bottom left: Phase-
space projection (dimensions 0, 1) showing agents collapsing to a common point. Bottom
right: Final statistics confirming zero variance across agents and preserved gauge symmetry.
linear group GL(K), which includes compact subgroups such as SO(N ) and SU(N ) and
more.
E.1 Hardware and Computational Details
All experiments were performed on a Windows workstation with 64GB RAM, AMD Ryzen
9900x CPU, and an Nvidia RTX5090 GPU.
Appendix F. Symmetry Breaking Simulations
To investigate observation-driven symmetry breaking, we simulated 8 overlapping agents
on a 2-dimensional flat base manifold under periodic boundary conditions, using SO(3)
with the ℓ = 4 irrep (9-dimensional representation space). All agents were initialized with
identical models and γij = 0.
F.0.1 Vacuum State
Without observations, all agent beliefs converge to a shared rotationally invariant vacuum
state with identical norms (Figure 11), occupying the gauge orbit predicted by the sym-
metric vacuum theory.
73
Dennis
Figure 12: Observation-driven symmetry breaking: μcenter
Q tracking summary with obser-
vations. Top: Norm trajectories diverge as each agent specializes to a distinct belief mag-
nitude. Middle left: Symmetry measure Var(∥μ∥) grows and saturates, indicating strong
symmetry breaking. Middle right: Population-averaged norm increases as agents acquire
information from observations. Bottom left: Phase-space projection (dimensions 0, 1)
showing agents separating into distinct trajectories. Bottom right: Final statistics con-
firming high variance across agents and strong symmetry breaking.
F.0.2 Observation-Driven Symmetry Breaking
With randomly drawn Gaussian observations, symmetry breaks and agents flow toward
unique norms (Figure 12), with distinct feature directions emerging as specialized modes of
the previously symmetric space.
These simulations use a different gauge group (SO(3)) and base manifold dimension
(2D) than the language modeling experiments (GL(10), 0D). The purpose is to validate the
general theoretical prediction that observations break gauge symmetry (Section 3.9), not to
directly model the language modeling setting. Whether the specific dynamics of symmetry
breaking transfer quantitatively across gauge groups and base manifold dimensions remains
an open question.
F.1 Model-Channel Formalism
The most general formulation of the gauge-covariant variational free energy includes a model
channel alongside the belief channel. Each agent i maintains, in addition to its belief qi(ki)
and prior pi(ki), a model distribution si(mi) over generative model parameters and a model
hyper-prior ri(mi). The complete free energy is (writing the partially minimized form with
βij and γij at their optimal softmax values, as discussed in Section 3.5):
74
Attention and Transformers as Free Energy Minimization
Ffull[{qi}, {si}] = X
i
DKL(qi ∥ pi) + X
i
DKL(si ∥ ri)
+ X
i,j
βij DKL(qi ∥ Ωij qj ) + X
i,j
γij DKL(si ∥ ˜Ωij sj )
− Eq[log p(o |{ki}, {mi})].
(190)
The model prior DKL(si∥ri) regularizes agent i’s model toward a hyper-prior ri de-
termined from a higher, slower meta-level. The model alignment term γij DKL(si∥ ˜Ωij sj )
enforces meta-cognitive consensus: agents with high γij are driven to agree on their gener-
ative models, implementing distributed model learning.
The prior pi(ki | mi) depends on model parameters that are themselves uncertain under
si(mi), creating a hierarchical Bayesian structure:
pi(ki) =
Z
pi(ki | mi) si(mi) dmi. (191)
The free energy exhibits natural timescale separation: beliefs qi (fast subsystem) update
rapidly while models si (slow subsystem) evolve on a slower timescale with ηs ≪ ηq. In the
present work, we set γij = 0 and hold models fixed, operating entirely in the fast (belief)
subsystem. Standard transformers also operate in this regime: inference (qi updates) pro-
ceeds with frozen parameters during a forward pass, while model parameters update via
backpropagation between forward passes.
References
P.-A. Absil, R. Mahony, and R. Sepulchre. Optimization Algorithms on Matrix Manifolds.
Princeton University Press, 2008.
Shun-ichi Amari. Natural gradient works efficiently in learning. Neural Computation, 10
(2):251–276, 1998.
Shun-ichi Amari. Information Geometry and Its Applications. Springer, 2016.
James A. Anderson and Edward Rosenfeld. Neurocomputing: Foundations of Research.
MIT Press, 1988. Often cited as Anderson 1984 for earlier technical report editions.
John C. Baez and Javier P. Muniain. Gauge Fields, Knots and Gravity. World Scientific,
1994.
Dzmitry Bahdanau, Kyunghyun Cho, and Yoshua Bengio. Neural machine translation by
jointly learning to align and translate. arXiv preprint arXiv:1409.0473, 2014.
Shaojie Bai, J. Zico Kolter, and Vladlen Koltun. Deep equilibrium models. In Advances in
Neural Information Processing Systems, volume 32, 2019.
Iz Beltagy, Matthew E Peters, and Arman Cohan. Longformer: The long-document trans-
former. arXiv preprint arXiv:2004.05150, 2020.
75
Dennis
David M. Blei, Alp Kucukelbir, and Jon D. McAuliffe. Variational inference: A review for
statisticians. Journal of the American Statistical Association, 112(518):859–877, 2017.
Rafal Bogacz. A tutorial on the free-energy framework for modelling perception and learn-
ing. Journal of Mathematical Psychology, 76:198–211, 2017.
Silvere Bonnabel. Stochastic gradient descent on Riemannian manifolds. IEEE Transactions
on Automatic Control, 58(9):2217–2229, 2013.
Michael M. Bronstein, Joan Bruna, Taco Cohen, and Petar Veliˇckovi´c. Geometric deep
learning: Grids, groups, graphs, geodesics, and gauges. arXiv preprint arXiv:2104.13478,
2021.
Stanley F. Chen and Joshua Goodman. An empirical study of smoothing techniques for
language modeling. Technical Report TR-10-98, Harvard University, 1998.
Kevin Clark, Urvashi Khandelwal, Omer Levy, and Christopher D. Manning. What does
BERT look at? An analysis of BERT’s attention. In Proceedings of the 2019 ACL
Workshop BlackboxNLP, pages 276–286, 2019.
Zihang Dai, Zhilin Yang, Yiming Yang, Jaime Carbonell, Quoc V Le, and Ruslan Salakhut-
dinov. Transformer-XL: Attentive language models beyond a fixed-length context. In
Proceedings of the 57th Annual Meeting of the Association for Computational Linguis-
tics, pages 2978–2988, 2019.
Yann N. Dauphin, Angela Fan, Michael Auli, and David Grangier. Language modeling with
gated convolutional networks. In International Conference on Machine Learning, pages
933–941. PMLR, 2017.
Jacob Devlin, Ming-Wei Chang, Kenton Lee, and Kristina Toutanova. BERT: Pre-
training of deep bidirectional transformers for language understanding. arXiv preprint
arXiv:1810.04805, 2018.
Nelson Elhage, Neel Nanda, Catherine Olsson, Tom Henighan, Nicholas Joseph, Ben Mann,
Amanda Askell, Yuntao Bai, Anna Chen, Tom Conerly, et al. A mathematical framework
for transformer circuits. Transformer Circuits Thread, 2021.
Marc Finzi, Samuel Stanton, Pavel Ohana, and Andrew Gordon Wilson. Generalizing
convolutional neural networks for equivariance to Lie groups on arbitrary continuous
data. In International Conference on Machine Learning, pages 3165–3176. PMLR, 2020.
Jakob N. Foerster, Yannis M. Assael, Nando de Freitas, and Shimon Whiteson. Learn-
ing to communicate with deep multi-agent reinforcement learning. Advances in Neural
Information Processing Systems, 29, 2016.
Theodore Frankel. The Geometry of Physics: An Introduction. Cambridge University Press,
3rd edition, 2011.
Karl Friston. The free-energy principle: A unified brain theory? Nature Reviews Neuro-
science, 11(2):127–138, 2010.
76
Attention and Transformers as Free Energy Minimization
Karl J. Friston, Thomas Parr, and Bert de Vries. The graphical brain: Belief propagation
and active inference. Network Neuroscience, 1(4):381–414, 2017.
Fabian B. Fuchs, Daniel E. Worrall, Volker Fischer, and Max Welling. SE(3)-transformers:
3D roto-translation equivariant attention networks. In Advances in Neural Information
Processing Systems, volume 33, pages 1970–1981, 2020.
William Fulton and Joe Harris. Representation Theory: A First Course, volume 129 of
Graduate Texts in Mathematics. Springer, 1991.
Roc´ıo Garc´ıa-Mill´an and Gunnar Pruessner. Network renormalization. Nature Physics, 20:
1114–1118, 2024.
Edouard Grave, Armand Joulin, and Nicolas Usunier. Improving neural language models
with a continuous cache. arXiv preprint arXiv:1612.04426, 2017.
Dan Hendrycks and Kevin Gimpel. Gaussian error linear units (GELUs). arXiv preprint
arXiv:1606.08415, 2016.
Geoffrey Hinton. The forward-forward algorithm: Some preliminary investigations. arXiv
preprint arXiv:2212.13345, 2022.
Albert Q Jiang, Alexandre Sablayrolles, Arthur Mensch, Chris Bamford, Devendra Singh
Chaplot, Diego de las Casas, Florian Bressand, Gianna Lengyel, Guillaume Lample, Lu-
cile Saulnier, et al. Mistral 7b. arXiv preprint arXiv:2310.06825, 2023.
Angelos Katharopoulos, Apoorv Vyas, Nikolaos Pappas, and Fran¸cois Fleuret. Transform-
ers are RNNs: Fast autoregressive transformers with linear attention. In International
Conference on Machine Learning, pages 5156–5165. PMLR, 2020.
Risi Kondor and Shubhendu Trivedi. On the generalization of equivariance and convolution
in neural networks to the action of compact groups. In International Conference on
Machine Learning, pages 2747–2755. PMLR, 2018.
Solomon Kullback and Richard A. Leibler. On information and sufficiency. Annals of
Mathematical Statistics, 22(1):79–86, 1951.
Stephen Merity, Caiming Xiong, James Bradbury, and Richard Socher. Pointer sentinel
mixture models. arXiv preprint arXiv:1609.07843, 2017.
Beren Millidge, Anil Seth, and Christopher L Buckley. Predictive coding: a theoretical and
experimental review. arXiv preprint arXiv:2107.12979, 2021.
Mikio Nakahara. Geometry, topology and physics. CRC Press, 2003.
Thomas Parr, Giovanni Pezzulo, and Karl J Friston. Active inference: The free energy
principle in mind, brain, and behavior. MIT Press, 2022.
Ofir Press, Noah A Smith, and Mike Lewis. Train short, test long: Attention with linear
biases enables input length generalization. International Conference on Learning Repre-
sentations, 2022.
77
Dennis
Alec Radford, Jeffrey Wu, Rewon Child, David Luan, Dario Amodei, and Ilya Sutskever.
Language models are unsupervised multitask learners. OpenAI Blog, 1(8):9, 2019.
Colin Raffel, Noam Shazeer, Adam Roberts, Katherine Lee, Sharan Narang, Michael
Matena, Yanqi Zhou, Wei Li, and Peter J Liu. Exploring the limits of transfer learning
with a unified text-to-text transformer. Journal of Machine Learning Research, 21(140):
1–67, 2020.
Prajit Ramachandran, Barret Zoph, and Quoc V. Le. Searching for activation functions.
arXiv preprint arXiv:1710.05941, 2018.
Hubert Ramsauer, Bernhard Sch¨afl, Johannes Lehner, Philipp Seidl, Michael Widrich,
Thomas Adler, Lukas Gruber, Markus Holzleitner, Milena Pavlovi´c, Geir Kjetil Sandve,
et al. Hopfield networks is all you need. In International Conference on Learning Repre-
sentations, 2021.
Maxwell J. D. Ramstead, Michael D. Kirchhoff, and Karl J. Friston. A tale of two densities:
Active inference is enactive inference. Adaptive Behavior, 28(4):225–239, 2020.
Rajesh PN Rao and Dana H Ballard. Predictive coding in the visual cortex: a functional
interpretation of some extra-classical receptive-field effects. Nature Neuroscience, 2(1):
79–87, 1999.
Xiao Shen, Hao Sun, and James P. LeSage. Coarse-graining as a route to microscopic
physics: The renormalization group in quantum field theory. Philosophical Transactions
of the Royal Society A, 369:1740–1758, 2011. Used in context of multi-scale agent coupling.
Ravid Shwartz-Ziv and Naftali Tishby. Opening the black box of deep neural networks via
information. arXiv preprint arXiv:1703.00810, 2017.
Shlomo Sternberg. Group Theory and Physics. Cambridge University Press, 1994.
Jianlin Su, Murtadha Ahmed, Yu Lu, Shengfeng Pan, Wen Bo, and Yunfeng Liu. RoFormer:
Enhanced transformer with rotary position embedding. Neurocomputing, 568:127063,
2024.
Nathaniel Thomas, Tess Smidt, Steven Kearnes, Lusann Yang, Li Li, Kai Kober, and Patrick
Riley. Tensor field networks: Rotation- and translation-equivariant neural networks for
3D point clouds. In arXiv preprint arXiv:1802.08219, 2018.
Naftali Tishby and Noga Zaslavsky. Deep learning and the information bottleneck principle.
IEEE Information Theory Workshop, pages 1–5, 2015.
Yao-Hung Hubert Tsai, Shaojie Bai, Rajan Puri, J Zico Kolter, Louis-Philippe Morency,
and Ruslan Salakhutdinov. Transformer dissection: An unified understanding for trans-
former’s attention via the lens of kernel. In Proceedings of the 2019 Conference on Em-
pirical Methods in Natural Language Processing, pages 4344–4353, 2019.
Ashish Vaswani, Noam Shazeer, Niki Parmar, Jakob Uszkoreit, Llion Jones, Aidan N.
Gomez, Lukasz Kaiser, and Illia Polosukhin. Attention is all you need. In Advances in
Neural Information Processing Systems, volume 30, pages 5998–6008, 2017.
78
Attention and Transformers as Free Energy Minimization
Maurice Weiler, Mario Geiger, Max Welling, Wouter Boomsma, and Taco Cohen. 3D steer-
able CNNs: Learning rotationally equivariant features in volumetric data. In Advances
in Neural Information Processing Systems, volume 31, 2018.
Steven Weinberg. The Quantum Theory of Fields, volume 2. Cambridge University Press,
1996. Volume 2: Modern Applications.
Kenneth G. Wilson and John Kogut. The renormalization group and the ϵ expansion.
Physics Reports, 12(2):75–199, 1974.
Michael Wooldridge. An Introduction to MultiAgent Systems. John Wiley & Sons, 2nd
edition, 2009.
79
Dennis
Table 1: Summary of notation. Symbols are grouped by category: base geometry, agents
and beliefs, gauge structure, variational quantities, and Transformer correspondences.
Symbol Domain Description
Base geometry and fibers
C Manifold Base manifold (spatial domain)
Bq (Eq ) RKq Belief fiber
Bp (Ep) RKp Model fiber
K (Kq ) Z+ Belief fiber (embedding) dimension
N Z+ Sequence length (number of agents)
Agents and beliefs
i, j {1, . . . , N } Agent (token) indices
qi = N (μq,i, Σq,i) P(RK ) Belief distribution of agent i (context-dependent)
pi = N (μp,i, Σp,i) P(RK ) Prior distribution (from token embedding)
si, ri P(RKp ) Model-fiber belief and prior (generative model)
Gauge geometry
G = GL(K) Lie group Structure group (general linear group)
g = gl(K) Lie algebra Lie algebra of G
ϕi ∈ g RK×K Gauge frame parameters for agent i (belief fiber)
˜ϕi ∈ g RKp×Kp Gauge frame parameters for agent i (model fiber)
Ωij = exp(ϕi) exp(−ϕj ) GL(K) Belief-fiber transport from frame j to frame i
˜Ωij GL(Kp) Model-fiber transport from frame j to frame i
Φi, ˜Φi RKp×Kq , RKq ×Kp Cross-fiber morphisms (belief ↔ model)
ρ : G → GL(V ) Representation of G on the fiber V
Variational quantities
F R Gauge-covariant variational free energy
DKL(q∥p) R≥0 Forward Kullback–Leibler divergence
βij [0, 1], P
j βij = 1 Belief attention weight (source-selection posterior)
γij [0, 1] Model-fiber coupling weight
α (αi) R+ Prior precision; state-dependent variant (Sec. 3.7)
κ (τ ) R+ Attention temperature (softmax inverse tempera-
ture)
λ R+ Belief alignment coupling strength
πj [0, 1], P
j πj = 1 Attention prior (encodes masking, positional bias)
Transformer correspondences (flat-bundle, isotropic limit)
Qi, Kj Rdk Query and key vectors (→ μi, μj under reduction)
Vj Rdv Value vector
WQ, WK Rdk ×d Projection matrices (→ gauge transformations)
dk Z+ Head dimension (= K per irrep block)
80
Attention and Transformers as Free Energy Minimization
Property Full Gauge Theory Standard Transformer
Gauge group G GL(K) or subgroup GL(dk) (implicit)
Gauge transformation Ωij = eϕi e−ϕj ∈ GL+(K) WQW ⊤
K = σ−2Ω−⊤ (learned,
constant)
Frame field ϕi(c) Varies per agent/position Absorbed into WQ, WK
Transport structure Agent-pair dependent Uniform across all pairs
Connection Aμ −∂μϕi + O(ϕ2) (0D: absent) 0 (flat connection)
Curvature / Holonomy Trivial (Ωij ΩjkΩki = I;
Sec. 6.3)
0 (trivial holonomy)
KL coupling DKL[qi ∥ Ωij qj ] DKL[qi ∥ Ωqj ] with constant Ω
What is learned Frame fields ϕi Global gauge WQ, WK
Table 2: Comparison between full gauge-theoretic formulation and standard transformers.
81
Dennis
VFE / Gauge Theory Limit / Mechanism Transformer Architecture Status
Attention Mechanism
βij = softmax(−DKL/τ ) Exact (Sec. 3.4) softmax(QK⊤/√dk) D
Ωij = eϕi e−ϕj ∈ GL+(K) Constant Ω (Limit 2) WQW ⊤
K = σ−2Ω−⊤ D
Forward KL uniqueness Thm. (App. C) Why dot-product attention D
mi = P
j βij Ωij μj Ω absorbed into WV
P
j αij Vj (value agg.) D
GL(dhead)H ⊂ GL(dk) Irrep decomposition Multi-head attention D
Positional Structure
πk = 0 for k > i Prior constraint Causal mask D
log πk = −m|i−k| Distance-decaying prior ALiBi D
Learnable πk(|i−k|) Learnable prior Relative pos. bias (T5) D
ϕ(pos) ∈ SO(2)dk /2 Pos.-dep. gauge frame RoPE D
πk = 0 for |i−k| > w Window prior Sliding window attention D
Normalization and Scaling
τ = √dk (σ2 absorbed) Concentration of measure 1/√dk scaling D
∥μj ∥2 ≈ C condition Gauge-geometric necessity Layer normalization S
gl(K) = sl(K) ⊕ R Remove trace component Mean centering in LN S
Training and Optimization
F[{qi}] Deterministic limit Loss L(θ) D
−Eq [log p(o|k)] Point evaluation limit Cross-entropy / MSE loss D
α DKL(qi∥pi) α = const, deterministic L2 weight decay D
α∗
i = c0/(b0 + DKL) State-dependent precision Adaptive weight decay D
dμ/dt = −η Σ∇F Euler discretization Gradient descent (backprop) D
Chain rule through layers Layered composition Backpropagation D
Vacuum (no observations) Gauge-symmetric minimum Untrained network D
Observations break symmetry Explicit breaking Training / learning D
Architecture
μ(ℓ+1) = μ(ℓ) − ˜η Σ∇F Natural gradient flow Residual connection D
e · exp(−∥e∥2/τ )/Z (Eq. 102) Boltzmann GLU gate GELU/SiLU activation D
VFE iterations (T inner, L outer) Shared / fresh weights Layers / depth; DEQ at L = 1 D
Prior bank {(μc
p, Σc
p, ϕc)} Deterministic limit Embedding matrix Wembed S
p(o|μ) = Cat(softmax(W μ)) Observation model Output proj. + softmax D
Fast subsystem Ffast Frozen models Forward pass (inference) S
Slow subsystem Fslow Between-pass updates Weight training S
Regularization
Stochastic πj → mj πj Random source masking Attention dropout I
Noise in μ → inflated Σ Effective uncertainty Hidden dropout I
Fisher-Rao G−1∇F Diagonal approx. (heuristic) Adam optimizer I
Riemannian trust region Norm clipping on manifold Gradient clipping I
Table 3: Comprehensive correspondence between the gauge-theoretic VFE framework and
transformer architecture. Status column: D = derived from the variational principle
(mathematically exact or follows from explicit limits); S = structural correspondence (the
framework explains the component’s role but does not uniquely predict its specific form); I
= interpretive.
82
Attention and Transformers as Free Energy Minimization
Model Architecture Params Grand mean ¯r ± SE
ALBERT-base-v2 12L × 12H 11.7M 0.851 0.008
BERT-base-uncased 12L × 12H 109.5M 0.804 0.012
BERT-large-uncased 24L × 16H 335.1M 0.778 0.013
DistilBERT-base 6L × 12H 66.4M 0.746 0.012
RoBERTa-base 12L × 12H 124.6M 0.612 0.006
Table 4: Multi-model validation at τ = 19.0 across 20 passages. All models use dhead = 64.
The α–β correspondence holds across architectures, with ¯r > 0.6 for all five models and
¯r > 0.8 for ALBERT and BERT-base. RoBERTa’s lower correlation reflects its different
pretraining regime (no next-sentence prediction, dynamic masking) which may induce at-
tention patterns less aligned with pure distance-based structure.
N Mean ¯r ± Std Min r Max r
16 0.790 0.220 −0.028 0.998
32 0.798 0.193 0.036 0.999
64 0.795 0.191 0.121 0.997
128 0.773 0.207 0.109 0.999
256 0.747 0.226 0.057 1.000
512 0.714 0.238 0.044 0.999
Table 5: Sequence-length sensitivity at τ = 19.0. The mean correlation remains above 0.71
across all lengths. The mild degradation at longer sequences is consistent with the key-norm
bias analysis: longer sequences increase key-norm variance (more tokens → more heteroge-
neous norms), amplifying the residual bias that the flat-bundle approximation neglects.
Model Layers Params Val PPL Test PPL
Gauge VFE (GL(10), K = 90) 1 58.8M 78.7 74.9
Modified KN-5 (matched BPE) — — 135.9 134.8
Standard transformer 1 17.6M 55.4 65.0
Table 6: Language modeling on WikiText-103 (vocab 50,257; random-chance PPL ≈
50,000). The gauge VFE model uses KL-divergence attention with no MLPs, pointwise
activation functions, or learned attention projections (WQ, WK , WV ); only a linear output
projection to vocabulary logits is retained. The modified Kneser-Ney 5-gram baseline (Chen
and Goodman, 1998) is trained on the same WikiText-103 data with matched GPT-2 BPE
tokenization (50,257 vocab) to ensure commensurable perplexity evaluation; see text for de-
tails. The standard transformer baseline uses a conventional architecture with dot-product
attention and MLP layers.
83
